% tenkz-grid — L2a: the contraction grid.
%
%   \begin{tenkz}[keys]  cell & cell & ... \\ ...  \end{tenkz}
%   \tnpic[keys]{cells}
%
% Rows are layers, columns are sites.  Bonds between horizontally adjacent
% cells are implicit; physical legs follow the typed row declarations
% (rows={ket|bra|op|wire}); adjacent facing legs pair automatically.  The
% picture is a math atom: its wire axis sits on the math axis.
%
% Boundary doctrine (0.6): open virtual indices are DEFAULT and are
% drawn as protruding stubs — a matrix product must look like a matrix.
% A picture side has exactly ONE policy word from the side-policy
% alphabet, `west=`/`east=` taking open | none | trace | cup | cup={m}:
% `none` seals the side, `trace` on both sides closes each row back to
% itself (alias: `periodic`), and `cup[={m}]` ties ADJACENT open rows to
% EACH OTHER in pairs by a smooth bend, optionally carrying a
% fixed-point matrix on the bend — see the cup section below.
% `boundary=` is the all-sides shorthand; rows override per side via
% rows= suffixes (`open`, `none`, `west none`, `east open`, ...).
% Every picture emits its boundary signature (open-index counts) to the
% event stream: equation sides must have equal signatures.
%
% Phase-0 scope: everything needed by benchmarks B1-B8 plus MPO/PEPS stress
% cases; `\tnfuse` supports rows=2; `periodic` closes the top row above and
% the bottom row below (single row: below).

\ExplSyntaxOn

% ---------- argument variants used below ----------
\cs_generate_variant:Nn \prop_put:Nnn { Nen, Nne, Nee }
\cs_generate_variant:Nn \prop_get:NnN { NeN }
\prg_generate_conditional_variant:Nnn \prop_get:NnN { NeN } { T, F, TF }
\prg_generate_conditional_variant:Nnn \prop_if_in:Nn { Ne } { p, T, F, TF }
\prg_generate_conditional_variant:Nnn \tl_if_blank:n { V } { T, F, TF }
\prg_generate_conditional_variant:Nnn \tl_if_in:nn { Vn } { T, F, TF }

% ---------- picture state ----------
\int_new:N \l__tenkz_nrows_int
\int_new:N \l__tenkz_ncols_int
\int_new:N \l__tenkz_row_int
\int_new:N \l__tenkz_col_int
\prop_new:N \l__tenkz_kind_prop      % "r-c" -> tn|tnc|tnx|dots|ghost|fuse|fusec
\prop_new:N \l__tenkz_label_prop     % "r-c" -> label tokens
\prop_new:N \l__tenkz_opts_prop      % "r-c" -> cell option tokens
\prop_new:N \l__tenkz_owner_prop     % "r-c" -> owning column (wide glyphs)
\prop_new:N \l__tenkz_name_prop      % "r-c" -> tikz node name
\seq_new:N \l__tenkz_rowtypes_seq    % per row: wire|ket|bra|op
\seq_new:N \l__tenkz_rowmods_seq     % per row: comma list of fused/operator
\seq_new:N \l__tenkz_spans_seq       % {r}{c}{k}{variant}{label}
\seq_new:N \l__tenkz_cuts_seq        % {r}{c}{label}
\seq_new:N \l__tenkz_bondlabels_seq  % {math}{c}{c'}
% Boundary model.  An open virtual index is data: a matrix product with
% sealed ends depicts a scalar or vector, not a matrix.  The picture-level
% policy is `open' (protruding stubs at each wire row's ends), `none'
% (sealed, for genuinely closed objects), or `periodic' (trace closure).
% Rows opt out/in via the rows= suffix grammar (rows={op:open, ket:none}).
\tl_new:N \l__tenkz_boundary_tl
\tl_new:N \l__tenkz_wlabel_tl    % boundary index label at the west stub tip
\tl_new:N \l__tenkz_elabel_tl    % boundary index label at the east stub tip
\int_new:N \l__tenkz_bvw_int     % open-index counts for the event signature
\int_new:N \l__tenkz_bve_int
\int_new:N \l__tenkz_bpu_int
\int_new:N \l__tenkz_bpd_int
\tl_new:N \l__tenkz_align_tl
% Inter-layer cups (`close west` / `close east`).  A cup bends row r's end
% into row r+1's end on one side: the ket-to-bra virtual contraction of a
% sandwich object.  This is NOT the periodic trace — `periodic` closes a
% row to ITSELF, while a cup contracts two DIFFERENT rows' indices.
\bool_new:N \l__tenkz_closew_bool
\bool_new:N \l__tenkz_closee_bool
\tl_new:N \l__tenkz_closewmat_tl % on-bend matrix of the west cups (may be empty)
\tl_new:N \l__tenkz_closeemat_tl % on-bend matrix of the east cups (may be empty)
\seq_new:N \l__tenkz_cupw_seq    % west cup pairs {top row}{bottom row}
\seq_new:N \l__tenkz_cupe_seq    % east cup pairs
\prop_new:N \l__tenkz_cuprow_prop% "west-r"/"east-r" -> row r is cup-closed there
\prop_new:N \l__tenkz_nopair_prop% rows carrying the `nopair` rows= modifier
% wires=k glyphs (one atom spanning k wire rows, quantikz \gate[k]).  Each
% spanned (row, owner-col) cell is keyed here to the glyph's node name, so
% the anchor resolver can terminate that row's wire ON the glyph's edge at
% the row's own height (the node's west/east anchors sit at the glyph's
% vertical CENTRE, which is between the wires).
\prop_new:N \l__tenkz_tall_prop
\dim_new:N \l__tenkz_tallx_dim
% physical trace (trace=physical): Tr_phys closes each site's up leg into
% its own down leg, so those legs leave the boundary signature
\bool_new:N \l__tenkz_ptrace_bool
\bool_new:N \l__tenkz_ptraceon_bool
% pair trace (trace={(physical, cols)}): the per-column partial
% trace of a ket-bra interface.  Per traced column the wrap loop IS Tr over
% that site's physical space -- it contracts the ket's physical index with
% the bra's below it -- so the column contributes NO open physical index.
% The traced ket cells are recorded "r-c" at render time (mirroring the
% nopair rows) so the expandable leg predicates can read the table.
\prop_new:N \l__tenkz_trcell_prop
\clist_new:N \l__tenkz_ptcols_clist
% bond direction (design brief): picture default, per-row resolution, and
% the resolved mid-wire mark style of the current row
\tl_new:N \l__tenkz_bonddir_tl
\tl_new:N \l__tenkz_dir_tl
\tl_new:N \l__tenkz_dirmarks_tl
% \tnskip holes: pending-west-stub flag of the bond scan, and a scratch tl
% for partner-cell kinds at pairing interfaces
\bool_new:N \l__tenkz_hole_bool
\tl_new:N \l__tenkz_pkind_tl
% The FRAME (frame=).  ONE linear map routes every drawn coordinate
% from the logical frame (columns run east, rows descend) to the paper
% frame.  frame=flat is the identity; frame=vertical is
% (x, y) -> (-y, -x): columns DESCEND the page and rows read west-to-east,
% so virtual bonds run vertically and physical legs point west (logical
% up) / east (logical down) -- the rotated-chain genre (RMP 2011.12127
% Fig:MPDO-O_L, the vertical O_L word; sec2fig8's rotated transfer-matrix
% column).  The map is orthogonal and an involution, so it also converts
% paper-frame compass words back to logical ones (see \tenkz_compass:n).
\bool_new:N \l__tenkz_vertical_bool
\dim_new:N \l__tenkz_mx_dim          % frame-map output: paper-frame point
\dim_new:N \l__tenkz_my_dim
\dim_new:N \l__tenkz_qx_dim          % saved mapped point (segment tails)
\dim_new:N \l__tenkz_qy_dim
% periodic closure style: `racetrack` (the explicit return wire) or
% `hooks` (the truncated half-hook closure of the RMP O_L figure)
\tl_new:N \l__tenkz_perstyle_tl
% per-side boundary policy (west= / east=), layered over boundary=.
% Empty = inherit the picture policy.  Prepared/discarded indices are
% SIDES of the word, not rows: a circuit fragment discards its west
% (input) indices while keeping its east (output) indices open, so the
% policy must be addressable per side (hard11).
\tl_new:N \l__tenkz_westpol_tl
\tl_new:N \l__tenkz_eastpol_tl

% ---------- environment keys ----------
\pgfkeys{
  /tenkz/grid/.is~family,
  /tenkz/grid/rows/.code={\tenkz_set_rows:n{#1}},
  /tenkz/grid/physical/.is~choice,
  /tenkz/grid/physical/up/.code={\tenkz_set_rows:n{ket}},
  /tenkz/grid/physical/down/.code={\tenkz_set_rows:n{bra}},
  /tenkz/grid/physical/updown/.code={\tenkz_set_rows:n{op}},
  /tenkz/grid/physical/none/.code={\tenkz_set_rows:n{wire}},
  /tenkz/grid/sandwich/.code={\tenkz_set_rows:n{ket,bra}},
  /tenkz/grid/boundary/.is~choice,
  /tenkz/grid/boundary/open/.code={\tl_set:Nn \l__tenkz_boundary_tl {open}},
  /tenkz/grid/boundary/none/.code={\tl_set:Nn \l__tenkz_boundary_tl {none}},
  % boundary=periodic / bare `periodic`: aliases for `west=trace,
  % east=trace` — the side-policy alphabet is the one boundary language,
  % and a trace on both sides IS the periodic closure
  /tenkz/grid/boundary/periodic/.code=
      {\tl_set:Nn \l__tenkz_westpol_tl {trace}
       \tl_set:Nn \l__tenkz_eastpol_tl {trace}},
  /tenkz/grid/periodic/.code=
      {\tl_set:Nn \l__tenkz_westpol_tl {trace}
       \tl_set:Nn \l__tenkz_eastpol_tl {trace}},
  % frame = flat|vertical: the picture's frame (see the frame note at the
  % picture state).  Vertical rotates the WHOLE grammar 90 degrees:
  % columns descend, bonds run vertically, physical legs point west/east
  % per row type, \tndots renders \vdots, the trace closure becomes a
  % side racetrack, and the label auto-quadrant rotates with the frame.
  % Flat stays the default identity.
  /tenkz/grid/frame/.is~choice,
  /tenkz/grid/frame/flat/.code=
      {\bool_set_false:N \l__tenkz_vertical_bool},
  /tenkz/grid/frame/vertical/.code=
      {\bool_set_true:N \l__tenkz_vertical_bool},
  % chain axis = east|south: alias of frame= naming the direction the
  % chain RUNS instead of the frame (east = flat, south = vertical) --
  % the spelling hard05_v2 anticipated.
  /tenkz/grid/chain~axis/.is~choice,
  /tenkz/grid/chain~axis/east/.code=
      {\bool_set_false:N \l__tenkz_vertical_bool},
  /tenkz/grid/chain~axis/south/.code=
      {\bool_set_true:N \l__tenkz_vertical_bool},
  % trace style = racetrack|hooks: HOW the trace closure is drawn, and
  % nothing else — a how-key never sets the whether (that is `west=trace,
  % east=trace`, or its alias `periodic`).  `racetrack` (flat default) is
  % the explicit return wire around the word.  `hooks` (vertical default)
  % is the RMP O_L closure (arXiv:2011.12127 Fig:MPDO-O_L): at each end
  % of the traced row the wire runs one half stub out, turns a
  % near-semicircle toward the return side (radius = stub/2, so the
  % hook's reach IS the stub ratio), and truncates just short of the
  % anchor line -- the ring seen edge-on, return wire implied.  Each hook
  % is ONE path (arc, not elbows), so its rounding is uniform by
  % construction (the plane_sweep3 1a/1b defects: to[]-curves render too
  % flat, multi-join racetracks mix sharp and rounded corners).
  /tenkz/grid/trace~style/.is~choice,
  /tenkz/grid/trace~style/racetrack/.code=
      {\tl_set:Nn \l__tenkz_perstyle_tl {racetrack}},
  /tenkz/grid/trace~style/hooks/.code=
      {\tl_set:Nn \l__tenkz_perstyle_tl {hooks}},
  % west = / east = <policy>: ONE side, ONE policy word from the 0.6
  % alphabet — open | none | trace | cup | cup={m}.  `open` draws the
  % side's protruding stubs, `none` seals it, `trace` closes each traced
  % row back to ITSELF (both sides trace = the periodic closure),
  % `cup[={m}]` ties ADJACENT open rows to EACH OTHER in pairs by a
  % smooth bend, optionally carrying a fixed-point matrix on the bend
  % (the 0.6 spelling of `close west` / `close east`).  Resolution per
  % row and side: the rows= suffix wins (`west none`, `east open`,
  % `open`, `none`), then the side policy, then the picture policy
  % (boundary=).  A row whose west is `none` draws no west stub but
  % keeps its east stub, and the boundary signature counts only the
  % stubs actually drawn -- prepared/discarded indices are sides, not
  % rows (hard11: the brick-circuit fragment seals its west inputs and
  % leaves its east outputs open).
  /tenkz/grid/west/.code={\tenkz_side_policy:nn {west} {#1}},
  /tenkz/grid/east/.code={\tenkz_side_policy:nn {east} {#1}},
  % The cup, in full (west=cup / east=cup={matrix}): inter-layer
  % closure.  On that side adjacent open wire rows are tied to EACH
  % OTHER in pairs, top-down (rows 1-2, 3-4, ...), each pair by one
  % smooth 180-degree bend (a cup): the ket-to-bra virtual contraction
  % of sandwich objects — orthogonality conditions (sum_s A_L^s{}^dagger
  % A_L^s = 1, RMP orthoL/orthoR), reduced density matrices (RMP
  % sec2fig8), tangent-space brackets (RMP PTA).
  % Topology note: this is DIFFERENT from periodic — a trace closes a row
  % back to ITSELF, while a cup contracts two DIFFERENT rows' indices.  An
  % odd row out keeps its stub; a cupped row draws no stub and contributes
  % no open index to the boundary signature on that side.  The optional
  % cup= value places a fixed-point matrix (filled dot) ON the bend apex,
  % label outside the bend — the rho^L / rho^R capping convention (RMP
  % sec2fig8, fig9-1); a bare `west=cup` / `east=cup` is a plain cup
  % (identity).
  % layer sep = <factor> multiplies the layer pitch \tenkz@r@layersep for
  % this picture only.  Motivation: a rail pair whose open legs point INTO
  % the inter-layer gap (both-rails-down reduced density matrices, RMP
  % sec2fig8 genre) needs roughly double the default gap so the upper
  % rail's leg-tip labels keep their clearance band.
  /tenkz/grid/layer~sep/.code={\cs_set:Npe \tenkz@r@layersep
      { \fp_eval:n { \tenkz@r@layersep * (#1) } }},
  % trace = {<cellset>} / open = {<cellset>}: the contraction policy of
  % interface SEGMENTS, addressed by the one cell-set algebra every tier
  % shares (terms (i, c) | (i1-i2, c1-c2) combined with +/-, top-level
  % commas as unions, resolved by the lattice resolver).  The first slot
  % names an INTERFACE: pairing interfaces are numbered top-down per
  % adjacent facing pair, and the word `physical' names the word's
  % physical self-interface, its topology resolved from row-type data,
  % not spelling — in a ket-over-bra sandwich `trace={(physical, 4-5)}`
  % closes columns 4-5 by the wrap loop (per column the per-site partial
  % trace Tr_s, so the picture depicts the marginal rho_kept =
  % Tr_traced |psi><psi|, RMP 2011.12127 fig. 2), while the bare word
  % `trace=physical` closes EVERY op-row site's up leg into its own down
  % leg (Tr_phys(M) = sum_s M^{ss}, the ZCL-MPDO genre; single-row
  % pictures only for now -- see \tenkz_draw_phtrace:).  `open={(i,c)}`
  % is the per-segment `nopair': the segment's pairing is suppressed and
  % both faces keep open stubs that count in the boundary signature.
  % Traced legs are CLOSED indices: they leave the signature.  Values
  % contain literal commas, so they MUST be braced.
  /tenkz/grid/trace/.code={\tenkz_trace_key:n{#1}},
  /tenkz/grid/open/.code={\tenkz_open_key:n{#1}},
  % bond dir = right|left: default direction of every wire row's bonds.  An
  % outgoing arrow is the vector space V, an incoming one its dual V*, and
  % contraction joins out-to-in (the mpo.tex school); a barb rides each bond
  % mid-wire.  Row-level override: the rows= modifiers `dir right`/`dir left`.
  % Directed boundary signatures (splitting virtual-west into in/out so that
  % oppositely-directed open legs compare UNEQUAL) are deferred to the next
  % audit revision.
  /tenkz/grid/bond~dir/.is~choice,
  /tenkz/grid/bond~dir/right/.code={\tl_set:Nn \l__tenkz_bonddir_tl {right}},
  /tenkz/grid/bond~dir/left/.code={\tl_set:Nn \l__tenkz_bonddir_tl {left}},
  /tenkz/grid/west~label/.store~in=\l__tenkz_wlabel_tl,
  /tenkz/grid/east~label/.store~in=\l__tenkz_elabel_tl,
  /tenkz/grid/bond~label/.code={\tenkz_bond_label:n{#1}},
  /tenkz/grid/align/.store~in=\l__tenkz_align_tl,
  /tenkz/grid/pitch/.code={\pgfkeys{/tenkz/pitch=#1}},
  /tenkz/grid/compact/.code={\pgfkeys{/tenkz/compact}},
  /tenkz/grid/inline/.code={\pgfkeys{/tenkz/inline}},
  /tenkz/grid/tensor~style/.code={\pgfqkeys{/tenkz}{tensor~style=#1}},
  % species={...}: picture-scope species declaration (see tenkz-core)
  /tenkz/grid/species/.code={\pgfqkeys{/tenkz}{species={#1}}},
}

% ---------- the side-policy alphabet (west= / east=) ----------
% One enum, applied per picture side: open | none | trace | cup |
% cup={m}.  `trace` records the side policy word (both sides trace = the
% periodic closure; one side alone is an error in waiting and warns at
% draw time); `cup[={m}]` routes to the inter-layer cup machinery.
\msg_new:nnn { tenkz } { side-policy }
  {
    Unknown~side~policy~'#2'~at~#1=.~A~picture~side~takes~exactly~one~
    word~of~the~policy~alphabet:~open,~none,~trace,~cup~or~cup={m}.
  }
\cs_new_protected:Npn \tenkz_side_word:nn #1#2
  {
    \str_if_eq:nnTF {#1} {west}
      { \tl_set:Nn \l__tenkz_westpol_tl {#2} }
      { \tl_set:Nn \l__tenkz_eastpol_tl {#2} }
  }
\cs_new_protected:Npn \tenkz_side_cup:nn #1#2
  {
    \str_if_eq:nnTF {#1} {west}
      { \bool_set_true:N \l__tenkz_closew_bool
        \tl_set:Nn \l__tenkz_closewmat_tl {#2} }
      { \bool_set_true:N \l__tenkz_closee_bool
        \tl_set:Nn \l__tenkz_closeemat_tl {#2} }
  }
% split a `cup=$m$' value at its (catcode-12) equals sign
\cs_new_protected:Npn \tenkz_side_cupval:nw #1 cup=#2 \q_stop
  { \tenkz_side_cup:nn {#1} {#2} }
\cs_new_protected:Npn \tenkz_side_policy:nn #1#2
  {
    \str_case:nnF {#2}
      {
        {open}  { \tenkz_side_word:nn {#1} {open}  }
        {none}  { \tenkz_side_word:nn {#1} {none}  }
        {trace} { \tenkz_side_word:nn {#1} {trace} }
        {cup}   { \tenkz_side_cup:nn  {#1} {}      }
      }
      {
        \tl_if_in:nnTF {#2} {cup=}
          { \tenkz_side_cupval:nw {#1} #2 \q_stop }
          { \msg_error:nnne {tenkz}{side-policy} {#1} { \tl_to_str:n {#2} } }
      }
  }
\cs_generate_variant:Nn \msg_error:nnnn { nnne }

% ---------- contraction cell-sets (open= / trace=) ----------
% Shared-resolver plumbing: the value is str-normalized, the word
% `physical' maps to the sentinel interface 0 (so the resolver sees
% pure integer tuples), top-level commas rewrite to unions through the
% lattice scanner, and the lattice resolver produces "i-c" membership
% keys.  Keys at interface 0 feed the physical-word column lists; keys
% at interface i >= 1 feed the interface-segment specs, mapped to rows
% at render time (interface numbering needs the row-type table).
\prop_new:N \l__tenkz_trspec_prop     % "i-c" -> traced interface segment
\prop_new:N \l__tenkz_openspec_prop   % "i-c" -> opened interface segment
\clist_new:N \l__tenkz_opencols_clist % `physical'-word opened columns
\prop_new:N \l__tenkz_opencell_prop   % render-time: "r-c" (r = the
                                      % interface's upper row) -> open
\int_new:N \l__tenkz_iface_int
\tl_new:N \l__tenkz_csexpr_tl
\prop_new:N \l__tenkz_csres_prop
\msg_new:nnn { tenkz } { contraction-set }
  {
    #1=~takes~a~braced~cell-set~of~(interface,~columns)~terms~--~
    interfaces~numbered~top-down~per~adjacent~facing~pair,~or~the~word~
    `physical'~--~e.g.~trace={(physical,~4-5)}~or~open={(1,2)}.~
    trace=~also~takes~the~bare~word~`physical'~(the~op~row's~
    self-trace).~Got~'#2'.
  }
\cs_new_protected:Npn \tenkz_cellset:nNN #1#2#3
  {
    \tl_set:Ne \l__tenkz_csexpr_tl { \tl_to_str:n {#1} }
    \use:e
      {
        \tl_replace_all:Nnn \exp_not:N \l__tenkz_csexpr_tl
          { \tl_to_str:n {physical} } { 0 }
      }
    \tl_clear:N \l__tenkzl_tmp_tl
    \int_zero:N \l__tenkzl_depth_int
    \tl_map_inline:Nn \l__tenkz_csexpr_tl { \tenkzl_openscan:n {##1} }
    \tenkzl_resolve:VN \l__tenkzl_tmp_tl \l__tenkz_csres_prop
    \prop_map_inline:Nn \l__tenkz_csres_prop
      { \tenkz_cellset_one:wNN ##1 \q_stop #2 #3 }
  }
\cs_new_protected:Npn \tenkz_cellset_one:wNN #1 - #2 \q_stop #3#4
  {
    \int_compare:nNnTF {#1} = {0}
      { \clist_put_right:Nn #4 {#2} }
      { \prop_put:Nnn #3 {#1-#2} {} }
  }
\cs_new_protected:Npn \tenkz_trace_key:n #1
  {
    \tl_if_in:nnTF {#1} {(}
      { \tenkz_cellset:nNN {#1} \l__tenkz_trspec_prop
          \l__tenkz_ptcols_clist }
      {
        \str_if_eq:nnTF {#1} {physical}
          { \bool_set_true:N \l__tenkz_ptrace_bool }
          { \msg_error:nnne {tenkz}{contraction-set} {trace}
              { \tl_to_str:n {#1} } }
      }
  }
\cs_new_protected:Npn \tenkz_open_key:n #1
  {
    \tl_if_in:nnTF {#1} {(}
      { \tenkz_cellset:nNN {#1} \l__tenkz_openspec_prop
          \l__tenkz_opencols_clist }
      { \msg_error:nnne {tenkz}{contraction-set} {open}
          { \tl_to_str:n {#1} } }
  }

\cs_new_protected:Npn \tenkz_set_rows:n #1
  {
    \seq_clear:N \l__tenkz_rowtypes_seq
    \seq_clear:N \l__tenkz_rowmods_seq
    \clist_map_inline:nn {#1}
      {
        % str-normalize so the document's catcode-12 colon matches
        \exp_args:Ne \tenkz_row_entry:n { \tl_to_str:n {##1} }
      }
  }
% split an entry  type[:mod[:mod]]  at the first (catcode-12) colon
\use:e
  {
    \cs_new_protected:Npn \exp_not:N \tenkz_row_entry_aux:w
      #1 \c_colon_str #2 \exp_not:N \q_stop
      {
        \seq_put_right:Nn \exp_not:N \l__tenkz_rowtypes_seq {#1}
        \seq_put_right:Nn \exp_not:N \l__tenkz_rowmods_seq {#2}
      }
  }
\cs_new_protected:Npn \tenkz_row_entry:n #1
  { \use:e { \exp_not:N \tenkz_row_entry_aux:w #1 \c_colon_str } \q_stop }

% `bond label = {$D$ at 1-2}` or `bond label = {$D$}` (defaults to bond 1-2)
\cs_new_protected:Npn \tenkz_bond_label:n #1
  { \tl_if_in:nnTF {#1} { at } { \tenkz_bond_label:w #1 \q_stop }
      { \seq_put_right:Nn \l__tenkz_bondlabels_seq { {#1}{1}{2} } } }
\cs_new_protected:Npn \tenkz_bond_label:w #1 at #2 - #3 \q_stop
  { \seq_put_right:Nn \l__tenkz_bondlabels_seq { {#1}{#2}{#3} } }

% ---------- cell commands (record; draw later) ----------
\cs_new:Npn \tenkz_cell_key:
  { \int_use:N \l__tenkz_row_int - \int_use:N \l__tenkz_col_int }
\cs_new_protected:Npn \tenkz_record:nnn #1#2#3
  {
    \prop_put:Nen \l__tenkz_kind_prop  { \tenkz_cell_key: } {#1}
    \prop_put:Nen \l__tenkz_opts_prop  { \tenkz_cell_key: } {#2}
    \prop_put:Nen \l__tenkz_label_prop { \tenkz_cell_key: } {#3}
  }

\NewDocumentCommand \tn { s O{} m }
  { \IfBooleanTF {#1} { \tenkz_record:nnn {tnc}{#2}{#3} }
                      { \tenkz_record:nnn {tn}{#2}{#3} } }
\NewDocumentCommand \tnX { O{} m }
  { \tenkz_record:nnn {tnx}{#1}{#2} }
\NewDocumentCommand \tnfuse { O{} m }
  { \tenkz_record:nnn {fuse}{#1}{#2} }
\NewDocumentCommand \tndots { O{} }
  { \tenkz_record:nnn {dots}{#1}{} }
\NewDocumentCommand \tnghost { m }
  { \tenkz_record:nnn {ghost}{}{#1} }
% \tnskip -- a HOLE: the site is absent, its indices are not.  A hole is an
% uncontracted site, not invisible ink: the column's virtual wires stay OPEN
% on both sides of the gap (no bond bridges it), and a physically paired
% partner's leg stays open too -- the tangent-space summand with the
% identity at site i (hard10) deletes the TENSOR, never the indices.
% Contrast \tnghost, which routes the bond straight through and silently
% swallows a paired partner's leg (that pairing now warns).
\NewDocumentCommand \tnskip { }
  { \tenkz_record:nnn {skip}{}{} }
\NewDocumentCommand \tnspan { O{brace~above} m m }
  {
    \seq_put_right:Ne \l__tenkz_spans_seq
      { { \int_use:N \l__tenkz_row_int }{ \int_use:N \l__tenkz_col_int }
        { #2 }{ \exp_not:n{#1} }{ \exp_not:n{#3} } }
  }
\NewDocumentCommand \tncut { O{} m }
  {
    \seq_put_right:Ne \l__tenkz_cuts_seq
      { { \int_use:N \l__tenkz_row_int }{ \int_use:N \l__tenkz_col_int }
        { \exp_not:n{#2} } }
  }

% ---------- per-cell option keys ----------
\tl_new:N \l__tenkz_cshape_tl
\tl_new:N \l__tenkz_cup_tl
\tl_new:N \l__tenkz_cdown_tl
\tl_new:N \l__tenkz_clpos_tl
\tl_new:N \l__tenkz_clshift_tl
\int_new:N \l__tenkz_cwide_int
\int_new:N \l__tenkz_cwires_int
\tl_new:N \l__tenkz_clegsat_tl
\bool_new:N \l__tenkz_cnolegs_bool
\tl_new:N \l__tenkz_crole_tl    % cell role word (empty = none)
\tl_new:N \l__tenkz_cspec_tl    % cell species name (empty = none)
% \tnfuse options
\int_new:N \l__tenkz_frows_int
\tl_new:N \l__tenkz_fcombined_tl

\pgfkeys{
  /tenkz/cell/.is~family,
  /tenkz/cell/dot/.code={\tl_set:Nn \l__tenkz_cshape_tl {dot}},
  /tenkz/cell/box/.code={\tl_set:Nn \l__tenkz_cshape_tl {box}},
  /tenkz/cell/pill/.code={\tl_set:Nn \l__tenkz_cshape_tl {pill}},
  /tenkz/cell/mpo/.code={\tl_set:Nn \l__tenkz_cshape_tl {mpo}},
  /tenkz/cell/wide/.code={\int_set:Nn \l__tenkz_cwide_int {#1}},
  % wires = <k>: the glyph spans k ADJACENT wire rows as ONE atom (quantikz
  % \gate[k]) and TERMINATES each spanned row's wire at its west and east
  % edges -- a two-site gate U acts on C^d (x) C^d, so both wires end ON the
  % gate and none passes through.  The skin grows to cover the rows; the
  % label stays centred.  \tnX takes the same key (a gauge capsule spanning
  % k rows of a doubled wire).
  /tenkz/cell/wires/.code={\int_set:Nn \l__tenkz_cwires_int {#1}},
  % legs at = {c1,c2,...}: a wide glyph's physical legs sit at the named
  % SPANNED columns (1-based within the span) instead of one centred leg --
  % the RFP isometry U_{(i_1 i_2), j} keeps one output leg PER BLOCKED SITE
  % (hard08), so the boundary signature counts one open physical index per
  % named column.  The up=/down= label is read as a comma list matched
  % leg-by-leg.
  /tenkz/cell/legs~at/.store~in=\l__tenkz_clegsat_tl,
  % tri = l|r: canonical-form isometry skins.  tri=l is the left-canonical
  % A_L (flat side west = domain, apex east = codomain, and the isometry
  % condition sum_s A_L^s{}^dagger A_L^s = 1 reads left-to-right); tri=r
  % mirrors it for A_R.  Label: inside under the box policy (if it fits),
  % beside in a free quadrant under the dot policy.
  /tenkz/cell/tri/.is~choice,
  /tenkz/cell/tri/l/.code={\tl_set:Nn \l__tenkz_cshape_tl {tril}},
  /tenkz/cell/tri/r/.code={\tl_set:Nn \l__tenkz_cshape_tl {trir}},
  /tenkz/cell/up/.store~in=\l__tenkz_cup_tl,
  /tenkz/cell/down/.store~in=\l__tenkz_cdown_tl,
  /tenkz/cell/label~pos/.store~in=\l__tenkz_clpos_tl,
  /tenkz/cell/label~shift/.store~in=\l__tenkz_clshift_tl,
  /tenkz/cell/no~legs/.code={\bool_set_true:N \l__tenkz_cnolegs_bool},
  % role= / species=: the semantic hue interface (see tenkz-core).  A
  % role or species resolves through pure style indirection — one
  % derived style per role x glyph family, `<role> tensor` for beads,
  % `<role> outline` for inscribed-label skins — so a theme rebinds
  % every hue at one point; the atom's event carries both words.  A
  % species must be declared (\tnset{species={...}}); role wins when
  % both are given.
  /tenkz/cell/role/.is~choice,
  /tenkz/cell/role/operator/.code={\tl_set:Nn \l__tenkz_crole_tl {operator}},
  /tenkz/cell/role/marked/.code={\tl_set:Nn \l__tenkz_crole_tl {marked}},
  /tenkz/cell/role/extra/.code={\tl_set:Nn \l__tenkz_crole_tl {extra}},
  /tenkz/cell/role/passive/.code={\tl_set:Nn \l__tenkz_crole_tl {passive}},
  /tenkz/cell/species/.code={\tenkz_species_check:n{#1}
      \tl_set:Nn \l__tenkz_cspec_tl {#1}},
  % \tnfuse: span=k is the 0.6 spelling (rows= already means row typing
  % at picture scope and would be its third meaning here); rows= stays
  % as the documented alias
  /tenkz/cell/span/.code={\int_set:Nn \l__tenkz_frows_int {#1}},
  /tenkz/cell/rows/.code={\int_set:Nn \l__tenkz_frows_int {#1}},
  /tenkz/cell/combined/.store~in=\l__tenkz_fcombined_tl,
}

\cs_new_protected:Npn \tenkz_cell_opts:n #1
  {
    % the default skin is the glyph policy (tensor style=dot|box), resolved
    % here at option-parse time; explicit per-cell skins override it below
    \tl_set:Ne \l__tenkz_cshape_tl { \tenkz@tensorstyle }
    \tl_clear:N \l__tenkz_cup_tl   \tl_clear:N \l__tenkz_cdown_tl
    \tl_clear:N \l__tenkz_clpos_tl \tl_clear:N \l__tenkz_clshift_tl
    \int_set:Nn \l__tenkz_cwide_int {1}
    \int_set:Nn \l__tenkz_cwires_int {1}
    \tl_clear:N \l__tenkz_clegsat_tl
    \bool_set_false:N \l__tenkz_cnolegs_bool
    \tl_clear:N \l__tenkz_crole_tl
    \tl_clear:N \l__tenkz_cspec_tl
    \int_set:Nn \l__tenkz_frows_int {2}
    \tl_set:Nn \l__tenkz_fcombined_tl {west}
    \tl_set:Nn \l__tenkz_gstrut_tl { \tenkz@glyphstrut }
    \tl_if_empty:nF {#1} { \pgfqkeys{/tenkz/cell}{#1} }
  }

% ---------- geometry ----------
\dim_new:N \l__tenkz_x_dim
\dim_new:N \l__tenkz_y_dim
% NB: eTeX dimen expressions want the integer factor on the RIGHT of `*`.
\cs_new:Npn \tenkz_colx:n #1
  { \dim_eval:n { \tenkz@pitch * (#1) } }
\cs_new:Npn \tenkz_rowy:n #1
  { \dim_eval:n { -\tenkz@r@layersep\tenkz@pitch * (#1-1) } }

% ---------- the frame map (frame=) ----------
% logical (x, y) -> paper (\l__tenkz_mx_dim, \l__tenkz_my_dim).
% horizontal: identity.  vertical: (x, y) -> (-y, -x), i.e. row 1's wire
% is the westmost column of the page and column 1 is the topmost site.
% The map is linear, orthogonal, and an involution (its own inverse), so
% straight bonds stay straight, circular arcs stay circular (only their
% start/end angles transform, see \tenkz_arcangle:n), and one function
% converts compass words in BOTH directions.
\cs_new_protected:Npn \tenkz_map:nn #1#2
  {
    \bool_if:NTF \l__tenkz_vertical_bool
      {
        % 0pt anchors the subtraction: eTeX rejects a unary minus
        % directly before a parenthesized subexpression
        \dim_set:Nn \l__tenkz_mx_dim { 0pt - ( #2 ) }
        \dim_set:Nn \l__tenkz_my_dim { 0pt - ( #1 ) }
      }
      {
        \dim_set:Nn \l__tenkz_mx_dim { #1 }
        \dim_set:Nn \l__tenkz_my_dim { #2 }
      }
  }
% the mapped point as a tikz coordinate (uses the LAST \tenkz_map:nn call)
\cs_new:Npn \tenkz_mpt: { ( \dim_use:N \l__tenkz_mx_dim , \dim_use:N \l__tenkz_my_dim ) }
% same, saved to (qx, qy) so a second map call can place a segment's far end
\cs_new_protected:Npn \tenkz_map_save:nn #1#2
  {
    \tenkz_map:nn {#1} {#2}
    \dim_set_eq:NN \l__tenkz_qx_dim \l__tenkz_mx_dim
    \dim_set_eq:NN \l__tenkz_qy_dim \l__tenkz_my_dim
  }
\cs_new:Npn \tenkz_qpt: { ( \dim_use:N \l__tenkz_qx_dim , \dim_use:N \l__tenkz_qy_dim ) }
% one logical-frame segment in style #5, both endpoints frame-mapped
\cs_new_protected:Npn \tenkz_seg:nnnnn #1#2#3#4#5
  {
    \tenkz_map_save:nn {#1} {#2}
    \tenkz_map:nn {#3} {#4}
    \draw[#5] \tenkz_qpt: -- \tenkz_mpt: ;
  }
% compass words under the frame map: the vertical map reflects across the
% north-west/south-east diagonal, so north<->west, south<->east, and the
% NW/SE diagonals are fixed.  Expandable (used inside anchor names); an
% involution, so it also converts a PAPER compass word back to a logical
% one (needed when a user-given `label pos` meets a logical-frame test).
\cs_new:Npn \tenkz_compass:n #1
  {
    \bool_if:NTF \l__tenkz_vertical_bool
      {
        \str_case:nn {#1}
          {
            {north}      {west}
            {south}      {east}
            {west}       {north}
            {east}       {south}
            {north~west} {north~west}
            {south~east} {south~east}
            {north~east} {south~west}
            {south~west} {north~east}
          }
      }
      {#1}
  }
\cs_generate_variant:Nn \tenkz_compass:n { V }
% arc angles under the frame map: a direction angle t maps to 270 - t
% (the reflection (x,y) -> (-y,-x) in angle form); sweep direction flips
% with it automatically because start and end both transform.
\cs_new:Npn \tenkz_arcangle:n #1
  {
    \bool_if:NTF \l__tenkz_vertical_bool
      { \int_eval:n { 270 - ( #1 ) } }
      { \int_eval:n {#1} }
  }
% the (dx,dy) paper offset of a logical west|east displacement of length
% #2 (a dimension expression): logical west maps to north (up the page),
% logical east to south -- (x, y) -> (-y, -x) sends (1, 0) to (0, -1).
% Expandable, so it can sit inside a tikz path (after ++) or a shift= key.
\cs_new:Npn \tenkz_side_vec:nn #1#2
  {
    \bool_if:NTF \l__tenkz_vertical_bool
      {
        \str_if_eq:nnTF {#1} {east}
          { (0pt,\dim_eval:n{0pt-(#2)}) }
          { (0pt,\dim_eval:n{#2}) }
      }
      {
        \str_if_eq:nnTF {#1} {east}
          { (\dim_eval:n{#2},0pt) }
          { (\dim_eval:n{0pt-(#2)},0pt) }
      }
  }
% the ++(dx,dy) offset of a west|east stub in the paper frame
\cs_new:Npn \tenkz_stub_vec:n #1
  { \tenkz_side_vec:nn {#1} { \tenkz@dim{\tenkz@r@stub} } }
% the (dx,dy) paper offset of a logical up|down physical displacement:
% outward sign #1 (1 = logical up, -1 = logical down), length #2 (a
% dimension expression).  Logical up maps to paper north (horizontal
% frame) or west (vertical frame: (x, y) -> (-y, -x) sends (0, 1) to
% (-1, 0)); the integer sign sits RIGHT of the `*` for eTeX.
\cs_new:Npn \tenkz_leg_vec:nn #1#2
  {
    \bool_if:NTF \l__tenkz_vertical_bool
      { (\dim_eval:n{ ( #2 ) * ( 0 - (#1) ) },0pt) }
      { (0pt,\dim_eval:n{ ( #2 ) * (#1) }) }
  }
% the paper direction of a logical up|down physical leg
\cs_new:Npn \tenkz_legdir:n #1
  {
    \bool_if:NTF \l__tenkz_vertical_bool
      { \str_if_eq:nnTF {#1} {up} {west} {east} }
      { \str_if_eq:nnTF {#1} {up} {north} {south} }
  }

% row-type queries (1-indexed); empty when out of range
\cs_new:Npn \tenkz_rowtype:n #1
  { \seq_item:Nn \l__tenkz_rowtypes_seq {#1} }
\cs_new:Npn \tenkz_rowmod:n #1
  { \seq_item:Nn \l__tenkz_rowmods_seq {#1} }

% get-or-clear: read a prop entry into a tl, emptying the tl on a miss so a
% failed lookup can never leak \q_no_value downstream (see the pitfall note
% on the anchor helper below).  The NeN variant takes a computed key.
\cs_new_protected:Npn \tenkz_get:NnN #1#2#3
  { \prop_get:NnNF #1 {#2} #3 { \tl_clear:N #3 } }
\cs_generate_variant:Nn \tenkz_get:NnN { NeN }

% nopair rows.  The rows= modifier `nopair` (rows={ket:nopair, bra})
% suppresses the automatic physical pairing across the interface at the
% modified row, so BOTH layers' legs stay open — the all-legs-open
% ket-over-bra window of the tangent-space projector (RMP PTA).  The
% pairing predicates are expandable and cannot run \str_if_in, so the
% renderer records the nopair rows in a prop at parse time (see
% \tenkz_render:nn) and the predicate reads that table.
\cs_new:Npn \tenkz_nopair_p:n #1
  { \prop_if_in_p:Ne \l__tenkz_nopair_prop { \int_eval:n {#1} } }

% traced cells (the trace= physical cell-set): expandable
% membership tests over the render-time table, mirroring the nopair
% mechanism.  \tenkz_traced_p:nn asks "does the ket cell (r,c) close by a
% wrap loop?"; \tenkz_traced_above_p:nn asks the bra-side question "is the
% cell just above (r,c) traced?" -- the loop consumes BOTH sites' legs.
\cs_new:Npn \tenkz_traced_p:nn #1#2
  { \prop_if_in_p:Ne \l__tenkz_trcell_prop
      { \int_eval:n {#1} - \int_eval:n {#2} } }
\cs_new:Npn \tenkz_traced_above_p:nn #1#2
  {
    \bool_lazy_and_p:nn
      { \int_compare_p:nNn {#1} > {1} }
      { \tenkz_traced_p:nn { #1 - 1 } {#2} }
  }

% pair[r] : row r's downward legs join row r+1's upward legs.
% An op row pairs onto anything below; a ket row pairs onto an op or bra
% below (the state's physical index feeds the operator or closes the
% sandwich).  Bra rows never provide downward legs.  A `nopair` modifier
% on either row of the interface suppresses the pairing (both legs open).
\cs_new:Npn \tenkz_pair_p:n #1
  {
    \bool_lazy_all_p:n
      {
        { ! \tenkz_nopair_p:n {#1} }
        { ! \tenkz_nopair_p:n {#1+1} }
        { \bool_lazy_or_p:nn
            { \bool_lazy_and_p:nn
                { \str_if_eq_p:ee { \tenkz_rowtype:n {#1} } {op} }
                { \bool_lazy_any_p:n
                    { { \str_if_eq_p:ee { \tenkz_rowtype:n {#1+1} } {op}  }
                      { \str_if_eq_p:ee { \tenkz_rowtype:n {#1+1} } {bra} }
                      { \str_if_eq_p:ee { \tenkz_rowtype:n {#1+1} } {ket} } } } }
            { \bool_lazy_and_p:nn
                { \str_if_eq_p:ee { \tenkz_rowtype:n {#1} } {ket} }
                { \bool_lazy_or_p:nn
                    { \str_if_eq_p:ee { \tenkz_rowtype:n {#1+1} } {bra} }
                    { \str_if_eq_p:ee { \tenkz_rowtype:n {#1+1} } {op} } } } }
      }
  }

% open legs of row r after pairing:  up / down -> true|false
% "up-paired" means row r-1 pairs down onto row r.
\cs_new:Npn \tenkz_uppaired_p:n #1
  {
    \bool_lazy_and_p:nn
      { \int_compare_p:nNn {#1} > {1} }
      { \tenkz_pair_p:n { #1 - 1 } }
  }
\cs_new:Npn \tenkz_openup_p:n #1
  {
    \bool_lazy_and_p:nn
      { ! \tenkz_uppaired_p:n {#1} }
      { \bool_lazy_or_p:nn
          { \str_if_eq_p:ee { \tenkz_rowtype:n {#1} } {op} }
          { \bool_lazy_and_p:nn
              { \str_if_eq_p:ee { \tenkz_rowtype:n {#1} } {ket} }
              { ! \tenkz_pair_p:n {#1} } } }
  }
\cs_new:Npn \tenkz_opendown_p:n #1
  {
    \bool_lazy_or_p:nn
      { \bool_lazy_and_p:nn
          { \str_if_eq_p:ee { \tenkz_rowtype:n {#1} } {op} }
          { ! \tenkz_pair_p:n {#1} } }
      { \bool_lazy_and_p:nn
          { \str_if_eq_p:ee { \tenkz_rowtype:n {#1} } {bra} }
          { ! \tenkz_uppaired_p:n {#1} } }
  }

% ---------- pair-trace recording ----------
% Populate the traced-cell table: a cell is traced when
% the trace= cell-set (interface-addressed, or the `physical' word)
% names its column in a ket row.  Eligibility is a ket-row tensor cell with
% a bra-row tensor cell directly below -- the per-site partial trace Tr_s
% contracts the ket's physical index with the bra's, so it needs both
% sites; anything else warns and stays untraced (its legs render as if no
% trace were asked).
\msg_new:nnn { tenkz } { pair-trace-cell }
  {
    the~traced~physical~column~at~row~#1,~column~#2~needs~a~tensor~cell~
    (\token_to_str:N \tn)~in~a~ket~row~with~a~tensor~cell~in~the~bra~row~
    directly~below:~the~per-site~partial~trace~contracts~the~ket's~
    physical~index~with~the~bra's.~The~cell~is~left~untraced.
  }
\cs_new_protected:Npn \tenkz_record_traces:
  {
    \prop_clear:N \l__tenkz_trcell_prop
    \int_step_inline:nn { \l__tenkz_nrows_int }
      {
        \int_step_inline:nnn {1} { \l__tenkz_ncols_int }
          {
            \bool_set_false:N \l_tmpa_bool
            % the physical cell-set: named columns of ket rows
            \str_if_eq:eeT { \tenkz_rowtype:n {##1} } {ket}
              {
                \clist_if_in:NnT \l__tenkz_ptcols_clist {####1}
                  { \bool_set_true:N \l_tmpa_bool }
              }
            \bool_if:NT \l_tmpa_bool
              { \tenkz_record_trace:nn {##1}{####1} }
          }
      }
  }
% eligibility check + record of one candidate (row #1, column #2)
\cs_new_protected:Npn \tenkz_record_trace:nn #1#2
  {
    \bool_set_false:N \l_tmpb_bool
    \str_if_eq:eeT { \tenkz_rowtype:n {#1} } {ket}
      {
        \str_if_eq:eeT { \tenkz_rowtype:n {#1+1} } {bra}
          {
            \tenkz_get:NeN \l__tenkz_kind_prop {#1-#2} \l__tenkz_pkind_tl
            \bool_lazy_or:nnT
              { \str_if_eq_p:Vn \l__tenkz_pkind_tl {tn} }
              { \str_if_eq_p:Vn \l__tenkz_pkind_tl {tnc} }
              {
                \tenkz_get:NeN \l__tenkz_kind_prop
                  { \int_eval:n {#1+1} - #2 } \l__tenkz_pkind_tl
                \bool_lazy_or:nnT
                  { \str_if_eq_p:Vn \l__tenkz_pkind_tl {tn} }
                  { \str_if_eq_p:Vn \l__tenkz_pkind_tl {tnc} }
                  { \bool_set_true:N \l_tmpb_bool }
              }
          }
      }
    \bool_if:NTF \l_tmpb_bool
      { \prop_put:Nen \l__tenkz_trcell_prop {#1-#2} {} }
      { \msg_warning:nnee {tenkz}{pair-trace-cell}
          { \int_eval:n {#1} }{ \int_eval:n {#2} } }
  }

% ---------- interface-addressed segment recording ----------
% Map the open=/trace= specs to rows: pairing interfaces are numbered
% top-down per adjacent facing pair (rows r with pair[r], in order), so
% interface i's segments live at (r_i, c) with r_i the i-th such row.
% Traced segments route through the pair-trace eligibility check (a
% non-ket-bra interface warns, as the per-cell key does); opened
% segments record directly — the draw pass guards with the leg
% candidacy.  The `physical' word's opened columns act on the pairing
% interface of a ket row, mirroring the traced-column sugar.
\msg_new:nnn { tenkz } { grid-interface-range }
  {
    open=/trace=~addresses~pairing~interface~#1,~but~this~picture~has~
    #2~pairing~interface(s),~numbered~top-down~per~adjacent~facing~
    pair:~no~segment~matches.
  }
\cs_new_protected:Npn \tenkz_record_interfaces:
  {
    \prop_clear:N \l__tenkz_opencell_prop
    \int_zero:N \l__tenkz_iface_int
    \int_compare:nNnT { \l__tenkz_nrows_int } > {1}
      {
        \int_step_inline:nn { \l__tenkz_nrows_int - 1 }
          {
            \bool_if:nT { \tenkz_pair_p:n {##1} }
              {
                \int_incr:N \l__tenkz_iface_int
                \int_step_inline:nnn {1} { \l__tenkz_ncols_int }
                  {
                    \prop_if_in:NeT \l__tenkz_trspec_prop
                      { \int_use:N \l__tenkz_iface_int - ####1 }
                      { \tenkz_record_trace:nn {##1}{####1} }
                    \prop_if_in:NeT \l__tenkz_openspec_prop
                      { \int_use:N \l__tenkz_iface_int - ####1 }
                      { \prop_put:Nen \l__tenkz_opencell_prop
                          {##1-####1} {} }
                    \str_if_eq:eeT { \tenkz_rowtype:n {##1} } {ket}
                      {
                        \clist_if_in:NnT \l__tenkz_opencols_clist {####1}
                          { \prop_put:Nen \l__tenkz_opencell_prop
                              {##1-####1} {} }
                      }
                  }
              }
          }
      }
    \prop_map_inline:Nn \l__tenkz_trspec_prop
      { \tenkz_iface_guard:w ##1 \q_stop }
    \prop_map_inline:Nn \l__tenkz_openspec_prop
      { \tenkz_iface_guard:w ##1 \q_stop }
  }
\cs_new_protected:Npn \tenkz_iface_guard:w #1 - #2 \q_stop
  {
    \int_compare:nNnT {#1} > { \l__tenkz_iface_int }
      { \msg_warning:nnee {tenkz}{grid-interface-range} {#1}
          { \int_use:N \l__tenkz_iface_int } }
  }

% ---------- the renderer ----------
\cs_new_protected:Npn \tenkz_render:nn #1#2
  {
    \group_begin:
    \prop_clear:N \l__tenkz_kind_prop
    \prop_clear:N \l__tenkz_label_prop
    \prop_clear:N \l__tenkz_opts_prop
    \prop_clear:N \l__tenkz_owner_prop
    \prop_clear:N \l__tenkz_name_prop
    \prop_clear:N \l__tenkz_tall_prop
    \bool_set_false:N \l__tenkz_ptrace_bool
    \clist_clear:N \l__tenkz_ptcols_clist
    \prop_clear:N \l__tenkz_trspec_prop
    \prop_clear:N \l__tenkz_openspec_prop
    \clist_clear:N \l__tenkz_opencols_clist
    \tl_clear:N \l__tenkz_bonddir_tl
    \seq_clear:N \l__tenkz_spans_seq
    \seq_clear:N \l__tenkz_cuts_seq
    \seq_clear:N \l__tenkz_bondlabels_seq
    \tl_set:Nn \l__tenkz_boundary_tl {open}
    \tl_clear:N \l__tenkz_wlabel_tl
    \tl_clear:N \l__tenkz_elabel_tl
    \tl_clear:N \l__tenkz_align_tl
    \bool_set_false:N \l__tenkz_closew_bool
    \bool_set_false:N \l__tenkz_closee_bool
    \tl_clear:N \l__tenkz_closewmat_tl
    \tl_clear:N \l__tenkz_closeemat_tl
    \seq_clear:N \l__tenkz_rowtypes_seq
    \seq_clear:N \l__tenkz_rowmods_seq
    \bool_set_false:N \l__tenkz_vertical_bool
    % empty = frame-dependent default, resolved at draw time (racetrack
    % for the horizontal frame, hooks for the vertical one)
    \tl_clear:N \l__tenkz_perstyle_tl
    \tl_clear:N \l__tenkz_westpol_tl
    \tl_clear:N \l__tenkz_eastpol_tl
    \tl_if_empty:nF {#1} { \pgfqkeys{/tenkz/grid}{#1} }
    % default row policy: bare wire rows
    \tenkz_parse_body:n {#2}
    % pad the row-type table with `wire`
    \int_while_do:nn
      { \seq_count:N \l__tenkz_rowtypes_seq < \l__tenkz_nrows_int }
      {
        \seq_put_right:Nn \l__tenkz_rowtypes_seq {wire}
        \seq_put_right:Nn \l__tenkz_rowmods_seq {}
      }
    % record the nopair rows for the expandable pairing predicates
    \prop_clear:N \l__tenkz_nopair_prop
    \int_step_inline:nn { \l__tenkz_nrows_int }
      {
        \tl_set:Ne \l_tmpb_tl { \tenkz_rowmod:n {##1} }
        \str_if_in:NnT \l_tmpb_tl {nopair}
          { \prop_put:Nen \l__tenkz_nopair_prop {##1} {} }
      }
    % record the row-scope side policies (rows={op:west none} etc.)
    \tenkz_record_rowsides:
    % record the traced ket cells for the pair-trace passes
    \tenkz_record_traces:
    % map the interface-addressed open=/trace= segments to rows
    \tenkz_record_interfaces:
    \tenkz@beginpicture{grid}
    \tenkz_draw:
    \group_end:
  }

\seq_new:N \l__tenkz_rows_seq
\seq_new:N \l__tenkz_cells_seq
\tl_new:N \l__tenkz_cell_tl

\cs_new_protected:Npn \tenkz_parse_body:n #1
  {
    \seq_set_split:Nnn \l__tenkz_rows_seq { \\ } {#1}
    % drop a trailing empty row produced by a final \\
    \seq_if_empty:NF \l__tenkz_rows_seq
      {
        \tl_set:Ne \l__tenkz_cell_tl { \seq_item:Nn \l__tenkz_rows_seq {-1} }
        \tl_if_blank:VT \l__tenkz_cell_tl
          { \seq_pop_right:NN \l__tenkz_rows_seq \l_tmpa_tl }
      }
    \int_set:Nn \l__tenkz_nrows_int { \seq_count:N \l__tenkz_rows_seq }
    \int_zero:N \l__tenkz_ncols_int
    \int_zero:N \l__tenkz_row_int
    \seq_map_inline:Nn \l__tenkz_rows_seq
      {
        \int_incr:N \l__tenkz_row_int
        \seq_set_split:Nnn \l__tenkz_cells_seq { & } {##1}
        \int_zero:N \l__tenkz_col_int
        \seq_map_inline:Nn \l__tenkz_cells_seq
          {
            \int_incr:N \l__tenkz_col_int
            \tl_set:Nn \l__tenkz_cell_tl {####1}
            \tl_if_blank:VF \l__tenkz_cell_tl { ####1 }
          }
        \int_set:Nn \l__tenkz_ncols_int
          { \int_max:nn { \l__tenkz_ncols_int } { \l__tenkz_col_int } }
      }
  }

% ---------- drawing ----------
\dim_new:N \l__tenkz_axis_dim
\dim_new:N \l__tenkz_tmp_dim

% frame=vertical covers the rotated-chain genre (plain chains and
% sandwiches: tensors, dots, bonds, stubs, legs, periodic closures, dot
% and boundary labels).  The remaining constructs still assume the
% horizontal frame; drawing them rotated would be silently wrong ink, so
% a vertical picture WARNS and skips them until the frame refactor
% reaches their passes.
\msg_new:nnn { tenkz } { vertical-unsupported }
  {
    frame=vertical~does~not~yet~transform~#1;~the~construct~is~
    skipped~in~this~picture.~Rotate~the~figure's~grammar~(plain~chains,~
    sandwiches,~periodic~closures~and~labels~are~covered)~or~keep~the~
    picture~horizontal.
  }
% a trace policy on ONE side only: a row closed back to itself needs
% both ends, so the closure cannot draw; the side stays sealed
\msg_new:nnn { tenkz } { trace-one-sided }
  {
    #1=trace~names~a~trace~closure~on~one~side~only;~a~row~closes~back~
    to~itself~through~BOTH~ends,~so~no~closure~is~drawn~(the~side~stays~
    sealed).~State~`west=trace,~east=trace'~--~or~its~alias~`periodic'.
  }

% drop the vertically-unsupported records, one warning each
\cs_new_protected:Npn \tenkz_vertical_prune:
  {
    \bool_lazy_or:nnT { \l__tenkz_closew_bool } { \l__tenkz_closee_bool }
      {
        \msg_warning:nnn {tenkz}{vertical-unsupported} {west=cup~/~east=cup}
        \bool_set_false:N \l__tenkz_closew_bool
        \bool_set_false:N \l__tenkz_closee_bool
      }
    \bool_if:NT \l__tenkz_ptrace_bool
      {
        \msg_warning:nnn {tenkz}{vertical-unsupported} {trace=physical}
        \bool_set_false:N \l__tenkz_ptrace_bool
      }
    \prop_if_empty:NF \l__tenkz_trcell_prop
      {
        \msg_warning:nnn {tenkz}{vertical-unsupported} {traced-column~wrap~loops}
        \prop_clear:N \l__tenkz_trcell_prop
      }
    \seq_if_empty:NF \l__tenkz_spans_seq
      {
        \msg_warning:nnn {tenkz}{vertical-unsupported} {\token_to_str:N \tnspan~braces}
        \seq_clear:N \l__tenkz_spans_seq
      }
    \seq_if_empty:NF \l__tenkz_cuts_seq
      {
        \msg_warning:nnn {tenkz}{vertical-unsupported} {\token_to_str:N \tncut~markers}
        \seq_clear:N \l__tenkz_cuts_seq
      }
    \seq_if_empty:NF \l__tenkz_bondlabels_seq
      {
        \msg_warning:nnn {tenkz}{vertical-unsupported} {bond~labels}
        \seq_clear:N \l__tenkz_bondlabels_seq
      }
  }

\cs_new_protected:Npn \tenkz_draw:
  {
    \bool_if:NT \l__tenkz_vertical_bool { \tenkz_vertical_prune: }
    % axis: chosen row, else the midline of the wire rows; a vertical
    % picture centres on the midline of its (descending) columns instead,
    % and align= (a ROW) has no vertical ordinate to name, so it is
    % ignored there
    \bool_if:NTF \l__tenkz_vertical_bool
      { \dim_set:Nn \l__tenkz_axis_dim
          { 0pt
            - ( \tenkz_colx:n {1} + \tenkz_colx:n {\l__tenkz_ncols_int} ) / 2 } }
      {
        \tl_if_empty:NTF \l__tenkz_align_tl
          { \dim_set:Nn \l__tenkz_axis_dim
              { ( \tenkz_rowy:n {1} + \tenkz_rowy:n {\l__tenkz_nrows_int} ) / 2 } }
          { \dim_set:Nn \l__tenkz_axis_dim
              { \tenkz_rowy:n {\l__tenkz_align_tl} } }
      }
    % NFSS sets math fonts up lazily; force them before reading the axis
    \check@mathfonts
    \dim_sub:Nn \l__tenkz_axis_dim { \fontdimen22\textfont2 }
    \begin{tikzpicture}
      [ tenkz~every~picture,
        baseline={(0pt,\dim_use:N \l__tenkz_axis_dim)} ]
      \tenkz_place_nodes:
      \tenkz_draw_bonds:
      \tenkz_compute_cups:
      \tenkz_draw_boundaries:
      \tenkz_draw_cups:
      \tenkz_draw_legs:
      \tenkz_draw_phtrace:
      \tenkz_draw_pairtrace:
      \tenkz_draw_dot_labels:
      \tenkz_draw_decorations:
      % the trace closure runs when BOTH side policies say trace (the
      % periodic alias sets exactly that); a one-sided trace has no
      % defined geometry — a row cannot close back to itself through one
      % end — so it warns and draws nothing
      \str_if_eq:VnTF \l__tenkz_westpol_tl {trace}
        {
          \str_if_eq:VnTF \l__tenkz_eastpol_tl {trace}
            {
              % closure-style default follows the frame: flat words keep
              % the explicit racetrack, vertical words default to hooks --
              % a vertical racetrack spends the very width the rotation
              % saves, so vertical trace closure must be compact by default
              \tl_if_empty:NT \l__tenkz_perstyle_tl
                {
                  \bool_if:NTF \l__tenkz_vertical_bool
                    { \tl_set:Nn \l__tenkz_perstyle_tl {hooks} }
                    { \tl_set:Nn \l__tenkz_perstyle_tl {racetrack} }
                }
              \str_if_eq:VnTF \l__tenkz_perstyle_tl {hooks}
                { \tenkz_draw_hooks: }
                { \tenkz_draw_trace: }
            }
            { \msg_warning:nnn {tenkz}{trace-one-sided} {west} }
        }
        {
          \str_if_eq:VnT \l__tenkz_eastpol_tl {trace}
            { \msg_warning:nnn {tenkz}{trace-one-sided} {east} }
        }
      \tenkz_emit_signature:
    \end{tikzpicture}
  }

% -- nodes ------------------------------------------------------------
\cs_new_protected:Npn \tenkz_place_nodes:
  {
    \int_step_inline:nn { \l__tenkz_nrows_int }
      {
        \int_step_inline:nnn {1} { \l__tenkz_ncols_int }
          { \tenkz_place_node:nn {##1} {####1} }
      }
  }

% the resolved role/species style of the current cell for glyph family
% #1 (tensor for beads, outline for inscribed-label skins), into
% \l__tenkz_rolesty_tl: species first, role wins — pure style
% indirection, no colour word ever appears at a call site.  The
% expandable event words report the resolved role/species (or `none').
\tl_new:N \l__tenkz_rolesty_tl
\cs_new_protected:Npn \tenkz_role_style:n #1
  {
    \tl_clear:N \l__tenkz_rolesty_tl
    % \c_space_tl, not ~: a space character after a control word is
    % swallowed at tokenization, so the style-name gap must be a token
    \tl_if_empty:NF \l__tenkz_cspec_tl
      { \tl_set:Ne \l__tenkz_rolesty_tl
          { species~ \tl_use:N \l__tenkz_cspec_tl \c_space_tl #1 } }
    \tl_if_empty:NF \l__tenkz_crole_tl
      { \tl_set:Ne \l__tenkz_rolesty_tl
          { \tl_use:N \l__tenkz_crole_tl \c_space_tl #1 } }
  }
% the resolved role style of the current SKIN: beads take the tensor
% family, every inscribed-label skin the outline family
\cs_new_protected:Npn \tenkz_role_style_skin:
  {
    \str_if_eq:VnTF \l__tenkz_cshape_tl {dot}
      { \tenkz_role_style:n {tensor} }
      { \tenkz_role_style:n {outline} }
  }
\cs_new:Npn \tenkz_role_word:
  { \tl_if_empty:NTF \l__tenkz_crole_tl {none}
      { \tl_use:N \l__tenkz_crole_tl } }
\cs_new:Npn \tenkz_spec_word:
  { \tl_if_empty:NTF \l__tenkz_cspec_tl {none}
      { \tl_use:N \l__tenkz_cspec_tl } }

% the current cell's kind, kept in a named local so it survives the
% \tenkz_cell_opts:n call and the label lookup below without a \l_tmpa_tl clash
\tl_new:N \l__tenkz_kind_tl
\cs_new_protected:Npn \tenkz_place_node:nn #1#2
  {
    \prop_get:NnNT \l__tenkz_kind_prop {#1-#2} \l__tenkz_kind_tl
      {
        \tenkz_get:NnN \l__tenkz_opts_prop {#1-#2} \l_tmpb_tl
        \exp_args:NV \tenkz_cell_opts:n \l_tmpb_tl
        \tenkz_get:NnN \l__tenkz_label_prop {#1-#2} \l_tmpb_tl
        \str_case:Vn \l__tenkz_kind_tl
          {
            {tn}   { \tenkz_node_tensor:nnV {#1}{#2} \l_tmpb_tl }
            {tnc}  { \tenkz_node_tensor:nnV {#1}{#2} \l_tmpb_tl }
            {tnx}  { \tenkz_node_simple:nnnV {#1}{#2}{on-wire~matrix} \l_tmpb_tl }
            {dots} { \tenkz_node_dots:nn {#1}{#2} }
            {ghost}{ \tenkz_node_ghost:nn {#1}{#2} }
            {fuse} { \tenkz_node_fuse:nn {#1}{#2} }
          }
      }
  }

\cs_new:Npn \tenkz_node_name:nn #1#2
  { tz \int_use:N \tenkz@pictureid - #1 - #2 }

% a tensor: dot (label outside) or box/pill/mpo (label inside)
\cs_new_protected:Npn \tenkz_node_tensor:nnn #1#2#3
  {
    % conjugate skin: overline the label
    \tenkz_get:NnN \l__tenkz_kind_prop {#1-#2} \l_tmpa_tl
    \str_if_eq:VnTF \l_tmpa_tl {tnc}
      { \tl_set:Nn \l__tenkz_cell_tl { \overline{#3} } }
      { \tl_set:Nn \l__tenkz_cell_tl { #3 } }
    % wide glyphs sit at the midpoint of their span and own every column
    \dim_set:Nn \l__tenkz_x_dim
      { \tenkz_colx:n {#2} + \tenkz@pitch * (\l__tenkz_cwide_int - 1) / 2 }
    \dim_set:Nn \l__tenkz_y_dim { \tenkz_rowy:n {#1} }
    \int_step_inline:nnn {#2} { #2 + \l__tenkz_cwide_int - 1 }
      { \prop_put:Nnn \l__tenkz_owner_prop {#1-##1} {#2} }
    % wires=k (quantikz \gate[k]): ONE atom centred between its outer wire
    % rows, TERMINATING each spanned row's wire at its edges.  Every spanned
    % cell is owned; continuation rows are marked `tnv` so nothing else
    % renders there; the tall table redirects each row's bonds and stubs to
    % the glyph's edge at that row's own height.  A bead cannot span wires:
    % the gate genre is a box, so a dot skin upgrades.
    \tenkz_tall_claim:nn {#1}{#2}
    \int_compare:nNnT { \l__tenkz_cwires_int } > {1}
      { \str_if_eq:VnT \l__tenkz_cshape_tl {dot}
          { \tl_set:Nn \l__tenkz_cshape_tl {box} } }
    % semantic hue: the resolved role/species style rides the xstyle
    % splice (labelled skins) or its own literal option slot (beads);
    % an empty slot is an inert trailing option
    \tenkz_role_style_skin:
    \tl_if_empty:NF \l__tenkz_rolesty_tl
      { \tl_put_right:Ne \l__tenkz_xstyle_tl
          { , \exp_not:V \l__tenkz_rolesty_tl } }
    % dot is the only skin without an inscribed label; box/mpo/pill share one
    % labelled-node path and differ only by their tikz style (pill also pins a
    % minimum width so a wide label grows the capsule sideways, not taller)
    \str_case:Vn \l__tenkz_cshape_tl
      {
        {dot}
          {
            \tenkz_map:nn { \l__tenkz_x_dim } { \l__tenkz_y_dim }
            \node[tensor, \l__tenkz_rolesty_tl]
              (\tenkz_node_name:nn{#1}{#2})
              at \tenkz_mpt: {};
          }
        {box}  { \tenkz_glyph_node:nnn {#1}{#2}{box~tensor} }
        {mpo}  { \tenkz_glyph_node:nnn {#1}{#2}{mpo~tensor} }
        {pill}
          { \tenkz_glyph_node:nnn {#1}{#2}
              { pill~tensor, minimum~width=\dim_eval:n
                  { \tenkz@pitch * (\l__tenkz_cwide_int - 1)
                    + 0.5\tenkz@pitch } } }
        % triangle skins inscribe the label only under the box policy; the
        % dot policy keeps the glyph empty and the dot-label pass places the
        % label beside it, in a free quadrant, like a bead's
        {tril} { \tenkz_tri_node:nnn {#1}{#2}{left~canonical~tensor} }
        {trir} { \tenkz_tri_node:nnn {#1}{#2}{right~canonical~tensor} }
      }
    \prop_put:Nne \l__tenkz_name_prop {#1-#2} { \tenkz_node_name:nn{#1}{#2} }
    \tenkz@event{atom|picture=\the\tenkz@pictureid|cell=#1-#2|kind=\tl_use:N \l__tenkz_cshape_tl|role=\tenkz_role_word:|species=\tenkz_spec_word:}
  }
\cs_generate_variant:Nn \tenkz_node_tensor:nnn { nnV }

% the ownership/geometry bookkeeping of a wires=k glyph at (#1,#2):
% recentre \l__tenkz_y_dim between the outer spanned rows, extend the extra
% node style by the covering minimum height (the glyph must overshoot its
% outer wires by \tenkz@r@tallover to read as terminating them), and claim
% every spanned cell.  No-op for k = 1.
\tl_new:N \l__tenkz_xstyle_tl
\cs_new_protected:Npn \tenkz_tall_claim:nn #1#2
  {
    \tl_clear:N \l__tenkz_xstyle_tl
    \int_compare:nNnT { \l__tenkz_cwires_int } > {1}
      {
        \dim_set:Nn \l__tenkz_y_dim
          { ( \tenkz_rowy:n {#1}
              + \tenkz_rowy:n { #1 + \l__tenkz_cwires_int - 1 } ) / 2 }
        % NB integer factor RIGHT of *; the layer pitch carries `layer sep`.
        % The covering dimension follows the frame: spanned rows stack
        % vertically in the horizontal frame (minimum height) and run
        % west-to-east in the vertical frame (minimum width).
        \tl_set:Ne \l__tenkz_xstyle_tl
          { , minimum~\bool_if:NTF \l__tenkz_vertical_bool {width}{height}
              =\dim_eval:n
              { \tenkz@r@layersep\tenkz@pitch * (\l__tenkz_cwires_int - 1)
                + \tenkz@dim{\tenkz@r@tallover} } }
        \int_step_inline:nnn { #1 } { #1 + \l__tenkz_cwires_int - 1 }
          {
            \prop_put:Nee \l__tenkz_tall_prop {##1-#2}
              { \tenkz_node_name:nn{#1}{#2} }
            \int_compare:nNnT {##1} > {#1}
              {
                \prop_put:Nen \l__tenkz_kind_prop {##1-#2} {tnv}
                \prop_put:Nee \l__tenkz_name_prop {##1-#2}
                  { \tenkz_node_name:nn{#1}{#2} }
                \int_step_inline:nnn {#2} { #2 + \l__tenkz_cwide_int - 1 }
                  { \prop_put:Nen \l__tenkz_owner_prop {##1-####1} {#2} }
              }
          }
      }
  }

% triangle node: inscribe the label only when the glyph policy inscribes
% (box); under the dot policy the label is external and the node is empty.
% Inscribed triangle labels drop the glyph strut: the strut equalizes BOX
% heights across a row, but a triangle must circumscribe its text box, so
% the strut's phantom super/subscripts would balloon the glyph until two
% triangle rows collide across the default layer gap.
\cs_new_protected:Npn \tenkz_tri_node:nnn #1#2#3
  {
    \str_if_eq:eeF { \tenkz@tensorstyle } {box}
      { \tl_clear:N \l__tenkz_cell_tl }
    \tl_clear:N \l__tenkz_gstrut_tl
    \tenkz_glyph_node:nnn {#1}{#2}{#3}
  }

% one labelled bead node: cell (#1,#2), tikz style list #3, drawn at the
% already-computed (\l__tenkz_x_dim,\l__tenkz_y_dim) with \l__tenkz_cell_tl;
% \l__tenkz_xstyle_tl carries the wires=k covering height (leading comma).
% The xstyle must be SPLICED into the option list before pgfkeys sees it:
% pgfkeys splits on top-level commas without expanding, so a comma hiding
% inside \tl_use:N would glue the height key onto the previous key's value.
\cs_new_protected:Npn \tenkz_glyph_node:nnn #1#2#3
  {
    \tl_set:Ne \l_tmpa_tl { \exp_not:n {#3} \exp_not:V \l__tenkz_xstyle_tl }
    \tenkz_glyph_node_aux:nnV {#1}{#2} \l_tmpa_tl
  }
\cs_new_protected:Npn \tenkz_glyph_node_aux:nnn #1#2#3
  {
    % (x, y) is the glyph centre in the LOGICAL frame; the frame map
    % places it (glyph shapes never rotate -- text stays upright)
    \tenkz_map:nn { \l__tenkz_x_dim } { \l__tenkz_y_dim }
    \node[#3]
      (\tenkz_node_name:nn{#1}{#2})
      at \tenkz_mpt:
      { $\tenkz@labelsize \tl_use:N \l__tenkz_gstrut_tl
        \tl_use:N \l__tenkz_cell_tl$ };
  }
\cs_generate_variant:Nn \tenkz_glyph_node_aux:nnn { nnV }
% the per-skin glyph strut: rectangular skins keep \tenkz@glyphstrut (boxes
% in one row share one height), triangle skins clear it (see above)
\tl_new:N \l__tenkz_gstrut_tl

% \tnX[wires=k] spans k rows of a doubled/stacked wire (a gauge transform
% acting on the fused index space, hard06/hard09 genre): same tall claim,
% the capsule's height grows to cover the rows
\cs_new_protected:Npn \tenkz_node_simple:nnnn #1#2#3#4
  {
    \dim_set:Nn \l__tenkz_x_dim { \tenkz_colx:n {#2} }
    \dim_set:Nn \l__tenkz_y_dim { \tenkz_rowy:n {#1} }
    \tenkz_tall_claim:nn {#1}{#2}
    % semantic hue: an on-wire matrix is an inscribed-label glyph, so a
    % role/species resolves through the outline family
    \tenkz_role_style:n {outline}
    \tl_if_empty:NF \l__tenkz_rolesty_tl
      { \tl_put_right:Ne \l__tenkz_xstyle_tl
          { , \exp_not:V \l__tenkz_rolesty_tl } }
    % splice the xstyle (see \tenkz_glyph_node:nnn on why it must be spliced)
    \tl_set:Nn \l__tenkz_cell_tl {#4}
    \tl_set:Ne \l_tmpa_tl { \exp_not:n {#3} \exp_not:V \l__tenkz_xstyle_tl }
    \tenkz_glyph_node_aux:nnV {#1}{#2} \l_tmpa_tl
    \prop_put:Nnn \l__tenkz_owner_prop {#1-#2} {#2}
    \prop_put:Nne \l__tenkz_name_prop {#1-#2} { \tenkz_node_name:nn{#1}{#2} }
    % an on-wire matrix is content ink: it must enter the audit's topology
    % hashes, or a gauge insertion X X^{-1} hashes equal to the bare chain
    \tenkz@event{atom|picture=\the\tenkz@pictureid|cell=#1-#2|kind=onwire|role=\tenkz_role_word:|species=\tenkz_spec_word:}
  }
\cs_generate_variant:Nn \tenkz_node_simple:nnnn { nnnV }

% the ellipsis follows the chain axis: \cdots along a horizontal word,
% \vdots down a vertical one (the O_L word's middle sites); the padding
% axis follows it, so the bond gaps around the dots stay symmetric
\cs_new_protected:Npn \tenkz_node_dots:nn #1#2
  {
    \tenkz_map:nn { \tenkz_colx:n {#2} } { \tenkz_rowy:n {#1} }
    \bool_if:NTF \l__tenkz_vertical_bool
      {
        \node[label, inner~ysep=0.30em]
          (\tenkz_node_name:nn{#1}{#2})
          at \tenkz_mpt: { $\vdots$ };
      }
      {
        \node[label, inner~xsep=0.30em]
          (\tenkz_node_name:nn{#1}{#2})
          at \tenkz_mpt: { $\cdots$ };
      }
    \prop_put:Nnn \l__tenkz_owner_prop {#1-#2} {#2}
    \prop_put:Nne \l__tenkz_name_prop {#1-#2} { \tenkz_node_name:nn{#1}{#2} }
  }

\cs_new_protected:Npn \tenkz_node_ghost:nn #1#2
  {
    \tenkz_map:nn { \tenkz_colx:n {#2} } { \tenkz_rowy:n {#1} }
    \node[inner~sep=0pt, minimum~size=0pt]
      (\tenkz_node_name:nn{#1}{#2})
      at \tenkz_mpt: {};
    \prop_put:Nnn \l__tenkz_owner_prop {#1-#2} {#2}
    \prop_put:Nne \l__tenkz_name_prop {#1-#2} { \tenkz_node_name:nn{#1}{#2} }
  }

% the fusion bar spans this row and the next (rows=2); its combined leg
% leaves west or east as a fused (doubled) stub
\cs_new_protected:Npn \tenkz_node_fuse:nn #1#2
  {
    \dim_set:Nn \l__tenkz_x_dim { \tenkz_colx:n {#2} }
    \dim_set:Nn \l__tenkz_y_dim { \tenkz_rowy:n {#1} }
    \dim_set:Nn \l__tenkz_tmp_dim
      { \tenkz_rowy:n { #1 + \l__tenkz_frows_int - 1 } }
    % the bar, with a small overhang beyond the outer wires (a logical
    % vertical stroke: the frame map turns it sideways with the rows)
    \tenkz_seg:nnnnn
      { \l__tenkz_x_dim } { \l__tenkz_y_dim + \tenkz@dim{0.10} }
      { \l__tenkz_x_dim } { \l__tenkz_tmp_dim - \tenkz@dim{0.10} }
      { fusion~map, line~width=\tenkz@dim{0.075} }
    % combined fused stub at the vertical midpoint
    \dim_set:Nn \l__tenkz_tmp_dim
      { ( \l__tenkz_y_dim + \l__tenkz_tmp_dim ) / 2 }
    % The combined leg is ONE index: V maps the pair of bond spaces into a
    % single fused bond space (V : C^{D_1} x C^{D_2} -> C^{D_12}), so it is
    % drawn as a single stroke.  A doubled stroke would misread as two
    % parallel indices.
    \str_if_eq:VnTF \l__tenkz_fcombined_tl {east}
      {
        \tenkz_seg:nnnnn
          { \l__tenkz_x_dim + \tenkz@dim{0.04} } { \l__tenkz_tmp_dim }
          { \l__tenkz_x_dim + \tenkz@dim{0.50} } { \l__tenkz_tmp_dim }
          { bond }
      }
      {
        \tenkz_seg:nnnnn
          { \l__tenkz_x_dim - \tenkz@dim{0.04} } { \l__tenkz_tmp_dim }
          { \l__tenkz_x_dim - \tenkz@dim{0.50} } { \l__tenkz_tmp_dim }
          { bond }
      }
    % the bar owns its column in every spanned row; continuation rows are
    % marked `fusec` so the node pass does not draw the bar twice
    \prop_put:Nnn \l__tenkz_owner_prop {#1-#2} {#2}
    \int_step_inline:nnn { #1 + 1 } { #1 + \l__tenkz_frows_int - 1 }
      {
        \prop_put:Nen \l__tenkz_kind_prop  {##1-#2} {fusec}
        \prop_put:Nen \l__tenkz_owner_prop {##1-#2} {#2}
      }
    % the bar label sits above the combined stub, clear of both wires
    \tenkz_get:NnN \l__tenkz_label_prop {#1-#2} \l_tmpb_tl
    \tl_if_blank:VF \l_tmpb_tl
      {
        \str_if_eq:VnTF \l__tenkz_fcombined_tl {east}
          { \dim_add:Nn \l__tenkz_x_dim { \tenkz@dim{0.30} } }
          { \dim_sub:Nn \l__tenkz_x_dim { \tenkz@dim{0.30} } }
        \tenkz_map:nn { \l__tenkz_x_dim }
          { \l__tenkz_tmp_dim + \tenkz@dim{\tenkz@r@labelclear} }
        \node[label, anchor=\tenkz_compass:n{south}] at \tenkz_mpt:
          { $\tenkz@labelsize \tl_use:N \l_tmpb_tl$ };
      }
    \tenkz@event{atom|picture=\the\tenkz@pictureid|cell=#1-#2|kind=fuse}
  }

% -- bonds ------------------------------------------------------------
% wire anchors of the item occupying (r,c):  west / east side
% (#1/#2 may be arbitrary integer expressions, e.g. an int register;
% \int_eval:n normalizes them before they become prop keys, and every
% \prop_get uses a T/F form so a miss can never leak \q_no_value)
\cs_new_protected:Npn \tenkz_anchor:nnnN #1#2#3#4
  {
    \prop_get:NeNF \l__tenkz_owner_prop
      {\int_eval:n{#1}-\int_eval:n{#2}} \l_tmpa_tl
      { \tl_set:Ne \l_tmpa_tl { \int_eval:n{#2} } }
    \tenkz_get:NeN \l__tenkz_kind_prop {\int_eval:n{#1}-\l_tmpa_tl} \l_tmpb_tl
    \bool_lazy_or:nnTF
      { \str_if_eq_p:Vn \l_tmpb_tl {fuse} }
      { \str_if_eq_p:Vn \l_tmpb_tl {fusec} }
      {
        % logical-frame point at the bar's edge, then frame-mapped
        \str_if_eq:nnTF {#3} {east}
          { \tenkz_map:nn
              { \tenkz_colx:n{\l_tmpa_tl} + \tenkz@dim{0.04} }
              { \tenkz_rowy:n {#1} } }
          { \tenkz_map:nn
              { \tenkz_colx:n{\l_tmpa_tl} - \tenkz@dim{0.04} }
              { \tenkz_rowy:n {#1} } }
        \tl_set:Ne #4 { \tenkz_mpt: }
      }
      {
        % a wires=k glyph: the node's west/east anchors sit at the glyph's
        % vertical CENTRE (between the wires), so each spanned row's wire
        % must instead end on the glyph's edge AT THE ROW'S OWN HEIGHT --
        % edge abscissa from the pgf shape anchor, ordinate from the row.
        % In the vertical frame the row's wire is a COLUMN of the page:
        % the edge ORDINATE comes from the (compass-mapped) shape anchor
        % and the abscissa is the row's mapped position.
        \prop_get:NeNTF \l__tenkz_tall_prop
          {\int_eval:n{#1}-\l_tmpa_tl} \l_tmpb_tl
          {
            \bool_if:NTF \l__tenkz_vertical_bool
              {
                \pgfextracty \l__tenkz_tallx_dim
                  { \exp_args:Ne \pgfpointanchor { \l_tmpb_tl }
                      { \tenkz_compass:n {#3} } }
                \tl_set:Ne #4
                  { ( \dim_eval:n { - \tenkz_rowy:n {#1} } ,
                      \dim_use:N \l__tenkz_tallx_dim ) }
              }
              {
                \pgfextractx \l__tenkz_tallx_dim
                  { \exp_args:Ne \pgfpointanchor { \l_tmpb_tl } {#3} }
                \tl_set:Ne #4
                  { ( \dim_use:N \l__tenkz_tallx_dim , \tenkz_rowy:n {#1} ) }
              }
          }
          {
            \tenkz_get:NeN \l__tenkz_name_prop
              {\int_eval:n{#1}-\l_tmpa_tl} \l_tmpb_tl
            \tl_set:Ne #4 { (\l_tmpb_tl. \tenkz_compass:n {#3}) }
          }
      }
  }

\tl_new:N \l__tenkz_pa_tl
\tl_new:N \l__tenkz_pb_tl

% extract a `species=<name>' row mod from the str-normalized mods #1
% (catcode-12 tokens, so both delimiters are built by \use:e): the name
% runs to the next colon; the appended sentinel provides both
% delimiters, so a species-free mods string yields an empty name.
\tl_new:N \l__tenkz_rowrole_tl
\tl_new:N \l__tenkz_rowspec_tl
\use:e
  {
    \cs_new_protected:Npn \exp_not:N \tenkz_row_species_aux:w
      #1 \tl_to_str:n {species=} #2 \c_colon_str #3 \exp_not:N \q_stop
      { \exp_not:N \tl_set:Nn \exp_not:N \l__tenkz_rowspec_tl {#2} }
  }
\cs_new_protected:Npn \tenkz_row_species:n #1
  {
    \use:e
      {
        \exp_not:N \tenkz_row_species_aux:w #1
        \c_colon_str \tl_to_str:n {species=} \c_colon_str
      } \q_stop
    \tl_remove_all:Nn \l__tenkz_rowspec_tl { ~ }
  }

% wire style of row #1 into #2.  Row mods were str-normalized at parse
% time (catcode-12 letters), so the membership tests must be string
% tests, not token tests.  The role words (`operator', `marked',
% `extra', `passive') and the `species=<name>' mod are the row-scope
% spellings of the semantic hue interface: each resolves to its derived
% bond style (style indirection, tenkz-core), the resolved words ride
% the row's bond events, and the closure passes (racetrack, hooks,
% cups) inherit the resolved style — an operator-hued traced row draws
% an operator-hued ring.
\cs_new_protected:Npn \tenkz_row_wirestyle:nN #1#2
  {
    \tl_set:Ne \l_tmpb_tl { \tenkz_rowmod:n {#1} }
    \tl_set:Nn #2 {bond}
    \tl_clear:N \l__tenkz_rowrole_tl
    \tl_clear:N \l__tenkz_rowspec_tl
    \str_if_in:NnT \l_tmpb_tl {fused}
      { \tl_set:Nn #2 {fused~bond} }
    \str_if_in:NnT \l_tmpb_tl {operator}
      { \tl_set:Nn #2 {operator~bond}
        \tl_set:Nn \l__tenkz_rowrole_tl {operator} }
    \str_if_in:NnT \l_tmpb_tl {marked}
      { \tl_set:Nn #2 {marked~bond}
        \tl_set:Nn \l__tenkz_rowrole_tl {marked} }
    \str_if_in:NnT \l_tmpb_tl {extra}
      { \tl_set:Nn #2 {extra~bond}
        \tl_set:Nn \l__tenkz_rowrole_tl {extra} }
    \str_if_in:NnT \l_tmpb_tl {passive}
      { \tl_set:Nn #2 {passive~bond}
        \tl_set:Nn \l__tenkz_rowrole_tl {passive} }
    \exp_args:NV \tenkz_row_species:n \l_tmpb_tl
    \tl_if_empty:NF \l__tenkz_rowspec_tl
      {
        \tenkz_species_check:V \l__tenkz_rowspec_tl
        \tl_set:Ne #2
          { species~ \tl_use:N \l__tenkz_rowspec_tl \c_space_tl bond }
      }
  }
% the row's resolved role/species event words (`none' when unhued)
\cs_new:Npn \tenkz_rowrole_word:
  { \tl_if_empty:NTF \l__tenkz_rowrole_tl {none}
      { \tl_use:N \l__tenkz_rowrole_tl } }
\cs_new:Npn \tenkz_rowspec_word:
  { \tl_if_empty:NTF \l__tenkz_rowspec_tl {none}
      { \tl_use:N \l__tenkz_rowspec_tl } }

% three-valued row direction (none|right|left): the rows= modifiers
% `dir right` / `dir left` override the picture-level `bond dir` default.
% Reversal is per-row mathematics -- the dual/inverse representation --
% so it flips the BARB (see the mark styles), never the wire.
\cs_new_protected:Npn \tenkz_row_dir:nN #1#2
  {
    \tl_set_eq:NN #2 \l__tenkz_bonddir_tl
    \tl_set:Ne \l_tmpb_tl { \tenkz_rowmod:n {#1} }
    \str_if_in:NnT \l_tmpb_tl {dir~right} { \tl_set:Nn #2 {right} }
    \str_if_in:NnT \l_tmpb_tl {dir~left}  { \tl_set:Nn #2 {left} }
  }
% the mid-wire mark style of the current row: barb forward for `right`
% (bonds are stroked west-to-east), reversed for `left`; the fused variant
% sizes the barb to span the doubled wire.  Empty when the row is
% undirected -- a trailing empty tikz option is inert.
\cs_new_protected:Npn \tenkz_dir_marks:NN #1#2
  {
    \tl_clear:N #2
    \tl_if_empty:NF \l__tenkz_dir_tl
      {
        \str_if_in:NnTF #1 {fused}
          { \tl_set:Nn #2 {fused~dir~marks} }
          { \tl_set:Nn #2 {dir~marks} }
        \str_if_eq:VnT \l__tenkz_dir_tl {left}
          { \tl_put_right:Nn #2 {~reversed} }
      }
  }

\cs_new_protected:Npn \tenkz_draw_bonds:
  {
    \int_step_inline:nn { \l__tenkz_nrows_int }
      {
        \tenkz_row_wirestyle:nN {##1} \l__tenkz_cell_tl
        \tenkz_row_dir:nN {##1} \l__tenkz_dir_tl
        \tenkz_dir_marks:NN \l__tenkz_cell_tl \l__tenkz_dirmarks_tl
        \bool_set_false:N \l__tenkz_hole_bool
        % connect consecutive occupied columns
        \int_zero:N \l_tmpa_int
        \int_step_inline:nnn {1} { \l__tenkz_ncols_int }
          {
            \tenkz_get:NeN \l__tenkz_kind_prop {##1-####1} \l__tenkz_pkind_tl
            \str_if_eq:VnTF \l__tenkz_pkind_tl {skip}
              {
                % a hole is an uncontracted site, not invisible ink: no bond
                % bridges the column, and the neighbours' virtual indices
                % facing the gap stay OPEN (protruding stubs) -- the
                % tangent-space summand with the identity at site i (hard10)
                % deletes the tensor, never the indices
                \int_compare:nNnT { \l_tmpa_int } > {0}
                  {
                    \tenkz_anchor:nnnN {##1}{\int_use:N\l_tmpa_int}
                      {east} \l__tenkz_pa_tl
                    \draw[\tl_use:N \l__tenkz_cell_tl]
                      \l__tenkz_pa_tl -- ++\tenkz_stub_vec:n{east};
                  }
                \int_zero:N \l_tmpa_int
                \bool_set_true:N \l__tenkz_hole_bool
                \tenkz@event{hole|picture=\the\tenkz@pictureid|row=##1|col=####1}
              }
              {
                \prop_if_in:NnT \l__tenkz_owner_prop {##1-####1}
                  {
                    \int_compare:nNnT { \l_tmpa_int } > {0}
                      {
                        \int_compare:nNnF { \l_tmpa_int } = { ####1 }
                          {
                            % same owner (a wide glyph) needs no bond
                            \tenkz_get:NeN \l__tenkz_owner_prop
                              {##1-\int_use:N\l_tmpa_int} \l_tmpa_tl
                            \tenkz_get:NeN \l__tenkz_owner_prop
                              {##1-####1} \l_tmpb_tl
                            \tl_if_eq:NNF \l_tmpa_tl \l_tmpb_tl
                              {
                                \tenkz_anchor:nnnN {##1}{\int_use:N\l_tmpa_int}
                                  {east} \l__tenkz_pa_tl
                                \tenkz_anchor:nnnN {##1}{####1}
                                  {west} \l__tenkz_pb_tl
                                \draw[\tl_use:N \l__tenkz_cell_tl,
                                      \tl_use:N \l__tenkz_dirmarks_tl]
                                  \l__tenkz_pa_tl -- \l__tenkz_pb_tl;
                                \tenkz@event{bond|picture=\the\tenkz@pictureid|row=##1|from=\int_use:N\l_tmpa_int|to=####1|dir=\tenkz_dir_word:|role=\tenkz_rowrole_word:|species=\tenkz_rowspec_word:}
                              }
                          }
                      }
                    % first tensor after a hole: its west index faces the
                    % gap and stays open
                    \bool_lazy_and:nnT
                      { \l__tenkz_hole_bool }
                      { \int_compare_p:nNn { \l_tmpa_int } = {0} }
                      {
                        \tenkz_anchor:nnnN {##1}{####1}{west} \l__tenkz_pb_tl
                        \draw[\tl_use:N \l__tenkz_cell_tl]
                          \l__tenkz_pb_tl -- ++\tenkz_stub_vec:n{west};
                        \bool_set_false:N \l__tenkz_hole_bool
                      }
                    \int_set:Nn \l_tmpa_int { ####1 }
                  }
              }
          }
      }
  }

% the event-stream word of the resolved row direction; from=/to= stay
% geometric, so `left` is recorded as `reverse`, not by swapping them
\cs_new:Npn \tenkz_dir_word:
  {
    \tl_if_empty:NTF \l__tenkz_dir_tl {none}
      { \str_if_eq:VnTF \l__tenkz_dir_tl {right} {forward} {reverse} }
  }

% -- physical legs ------------------------------------------------------
% Legs attach only to tensor glyphs (kind tn / tnc), never to on-wire
% matrices, dots, ghosts, or fusion bars.
\cs_new_protected:Npn \tenkz_leg_candidate:nnT #1#2#3
  {
    \prop_get:NeNT \l__tenkz_kind_prop {#1-#2} \l_tmpa_tl
      {
        \bool_lazy_or:nnT
          { \str_if_eq_p:Vn \l_tmpa_tl {tn} }
          { \str_if_eq_p:Vn \l_tmpa_tl {tnc} }
          {
            % respect `no legs`
            \tenkz_get:NeN \l__tenkz_opts_prop {#1-#2} \l_tmpb_tl
            \exp_args:NV \tenkz_cell_opts:n \l_tmpb_tl
            \bool_if:NF \l__tenkz_cnolegs_bool {#3}
          }
      }
  }

% pairing onto a \tnghost cell silently swallows the live partner's leg --
% invisible ink is a trap when the reader expects a contraction or an open
% index.  The warning names the fix: a hole is \tnskip, not \tnghost.
\msg_new:nnn { tenkz } { ghost-pair }
  {
    A~physical~pairing~meets~a~\token_to_str:N \tnghost \c_space_tl
    cell~at~row~#1,~column~#2:~the~partner's~leg~is~dropped.~A~ghost~is~
    invisible~ink,~not~a~hole~--~use~\token_to_str:N \tnskip \c_space_tl
    if~the~site~is~an~uncontracted~hole~(the~partner's~leg~then~stays~open).
  }
% message arguments are not expanded by l3msg, so row/column expressions
% must be e-expanded at the call site or the log prints \int_eval:n {...}
\cs_generate_variant:Nn \msg_warning:nnnn { nnee }

\cs_new_protected:Npn \tenkz_draw_legs:
  {
    \int_step_inline:nn { \l__tenkz_nrows_int }
      {
        \int_step_inline:nnn {1} { \l__tenkz_ncols_int }
          {
            \prop_get:NnNT \l__tenkz_name_prop {##1-####1} \l__tenkz_pa_tl
              {
                \tenkz_leg_candidate:nnT {##1}{####1}
                  {
                    % a traced cell's legs belong to its wrap loop (the
                    % per-site Tr_s draws them), so the open-stub passes
                    % skip it -- and skip the bra partner just below it
                    \bool_lazy_and:nnT
                      { \tenkz_openup_p:n {##1} }
                      { ! \tenkz_traced_p:nn {##1}{####1} }
                      { \tenkz_leg_open:nnn {##1}{####1}{up} }
                    \bool_lazy_and:nnT
                      { \tenkz_opendown_p:n {##1} }
                      { ! \tenkz_traced_above_p:nn {##1}{####1} }
                      { \tenkz_leg_open:nnn {##1}{####1}{down} }
                  }
                % paired leg to the row below
                \bool_if:nT { \tenkz_pair_p:n {##1} }
                  {
                    \tenkz_leg_candidate:nnT {##1}{####1}
                      {
                        \tenkz_get:NeN \l__tenkz_kind_prop
                          { \int_eval:n{##1+1} - ####1 } \l__tenkz_pkind_tl
                        \str_case:VnF \l__tenkz_pkind_tl
                          {
                            % a hole below the interface: the pairing has no
                            % partner, so THIS tensor's down leg stays open
                            % (an uncontracted physical index)
                            {skip}
                              { \tenkz_leg_open:nnn {##1}{####1}{down} }
                            {ghost}
                              { \msg_warning:nnee {tenkz}{ghost-pair}
                                  { \int_eval:n{##1+1} }{####1} }
                          }
                          {
                            \prop_get:NeNT \l__tenkz_name_prop
                              { \int_eval:n{##1+1} - ####1 } \l__tenkz_pb_tl
                              {
                                \tenkz_leg_candidate:nnT
                                  { \int_eval:n{##1+1} }{####1}
                                  {
                                    % a traced segment's wrap loop is
                                    % drawn by the pair-trace pass; an
                                    % OPENED segment (open={(i,c)}, the
                                    % per-segment nopair) keeps both
                                    % faces' stubs instead of a wire
                                    \bool_if:nF
                                      { \tenkz_traced_p:nn {##1}{####1} }
                                      {
                                        \prop_if_in:NeTF
                                          \l__tenkz_opencell_prop
                                          {##1-####1}
                                          {
                                            \tenkz_leg_open:nnn
                                              {##1}{####1}{down}
                                            \tenkz_leg_open:enn
                                              { \int_eval:n{##1+1} }
                                              {####1}{up}
                                          }
                                          {
                                            % logical south/north faces of
                                            % the interface; the compass
                                            % map turns them east/west in
                                            % a vertical frame
                                            \draw[physical~leg]
                                              (\l__tenkz_pa_tl.\tenkz_compass:n{south}) --
                                              (\l__tenkz_pb_tl.\tenkz_compass:n{north});
                                          }
                                      }
                                  }
                              }
                          }
                      }
                  }
              }
            % a hole (or ghost) ABOVE a paired interface: the live partner
            % below keeps its up leg open (hole), or loses it audibly (ghost)
            \int_compare:nNnT {##1} < { \l__tenkz_nrows_int }
              {
                \bool_if:nT { \tenkz_pair_p:n {##1} }
                  {
                    \tenkz_get:NeN \l__tenkz_kind_prop
                      {##1-####1} \l__tenkz_pkind_tl
                    \str_case:Vn \l__tenkz_pkind_tl
                      {
                        {skip}
                          { \tenkz_leg_candidate:nnT
                              { \int_eval:n{##1+1} }{####1}
                              { \tenkz_leg_open:enn
                                  { \int_eval:n{##1+1} }{####1}{up} } }
                        {ghost}
                          { \tenkz_leg_candidate:nnT
                              { \int_eval:n{##1+1} }{####1}
                              { \msg_warning:nnnn {tenkz}{ghost-pair}
                                  {##1}{####1} } }
                      }
                  }
              }
          }
      }
  }

% an open leg plus its tip label, placed clear of the leg by construction.
% Direction is data: `up` draws from the north anchor outward (sign +1) with
% the label below-anchored above the tip; `down` mirrors it from the south
% anchor (sign -1).  Only the anchors, the label slot, and the sign differ.
\cs_new_protected:Npn \tenkz_leg_open:nnn #1#2#3
  {
    \tenkz_get:NnN \l__tenkz_name_prop {#1-#2} \l__tenkz_pa_tl
    \tenkz_get:NnN \l__tenkz_opts_prop {#1-#2} \l_tmpb_tl
    \exp_args:NV \tenkz_cell_opts:n \l_tmpb_tl
    \tl_if_empty:NTF \l__tenkz_clegsat_tl
      {
        \str_if_eq:nnTF {#3} {up}
          { \tenkz_leg_stub:nnNn {north}{south} \l__tenkz_cup_tl   { 1} }
          { \tenkz_leg_stub:nnNn {south}{north} \l__tenkz_cdown_tl {-1} }
      }
      {
        % legs at={c1,...}: one leg PER NAMED SPANNED COLUMN of the wide
        % glyph -- each an independent open physical index (U_{(i1 i2),j}
        % keeps one output per blocked site, hard08)
        \str_if_eq:nnTF {#3} {up}
          { \tenkz_leg_multi:nnnnNn {#1}{#2}{north}{south}
              \l__tenkz_cup_tl   { 1} }
          { \tenkz_leg_multi:nnnnNn {#1}{#2}{south}{north}
              \l__tenkz_cdown_tl {-1} }
      }
  }
\cs_generate_variant:Nn \tenkz_leg_open:nnn { enn }

% the multi-leg form: row #1, base column #2, glyph-edge anchor #3, label
% anchor #4, label-list var #5 (a comma list matched leg-by-leg, popped
% without expansion so math tokens survive), outward sign #6.  The edge
% ordinate comes from the node's shape anchor (one rectangle edge, so it
% is the same at every spanned column); each leg's abscissa is its named
% column's -- the leg sits over the blocked site it indexes.
\seq_new:N \l__tenkz_leglab_seq
\tl_new:N \l__tenkz_legbase_tl
\cs_generate_variant:Nn \seq_set_from_clist:Nn { NV }
\cs_new_protected:Npn \tenkz_leg_multi:nnnnNn #1#2#3#4#5#6
  {
    % the glyph-edge coordinate is one rectangle edge, so ONE paper
    % ordinate serves every spanned column: the edge's y in the
    % horizontal frame, its x in the vertical one (where the edge anchor
    % word itself passes through the compass map)
    \bool_if:NTF \l__tenkz_vertical_bool
      {
        \pgfextractx \l__tenkz_tmp_dim
          { \exp_args:Ne \pgfpointanchor { \l__tenkz_pa_tl }
              { \tenkz_compass:n {#3} } }
      }
      {
        \pgfextracty \l__tenkz_tmp_dim
          { \exp_args:Ne \pgfpointanchor { \l__tenkz_pa_tl } {#3} }
      }
    \seq_set_from_clist:NV \l__tenkz_leglab_seq #5
    \exp_args:NV \clist_map_inline:nn \l__tenkz_clegsat_tl
      {
        % leg base = (column abscissa, edge ordinate) in the logical
        % frame; the vertical frame maps the column abscissa x to paper
        % y = -x and keeps the extracted edge coordinate as paper x
        \dim_set:Nn \l__tenkz_x_dim { \tenkz_colx:n { #2 + ##1 - 1 } }
        \bool_if:NTF \l__tenkz_vertical_bool
          { \tl_set:Ne \l__tenkz_legbase_tl
              { ( \dim_use:N \l__tenkz_tmp_dim ,
                  \dim_eval:n { -\l__tenkz_x_dim } ) } }
          { \tl_set:Ne \l__tenkz_legbase_tl
              { ( \dim_use:N \l__tenkz_x_dim ,
                  \dim_use:N \l__tenkz_tmp_dim ) } }
        % offsets as ++/shift vectors, never inside calc (see leg_stub)
        \draw[physical~leg]
          \l__tenkz_legbase_tl --
          ++\tenkz_leg_vec:nn{#6}{\tenkz@dim{\tenkz@r@physleg}};
        \seq_pop_left:NNT \l__tenkz_leglab_seq \l_tmpb_tl
          {
            \tl_if_blank:VF \l_tmpb_tl
              {
                \node[label, anchor=\tenkz_compass:n{#4},
                      shift={\tenkz_leg_vec:nn{#6}
                        { \tenkz@dim{\tenkz@r@physleg}
                          + \tenkz@dim{\tenkz@r@labelclear} }}]
                  at \l__tenkz_legbase_tl
                  { \tl_use:N \l_tmpb_tl };
              }
          }
      }
  }
% #1 LOGICAL node anchor, #2 logical label anchor, #3 label var, #4
% outward sign (1|-1).  Both anchors and the outward offset pass through
% the frame map: in a vertical picture an `up` leg leaves the glyph's
% paper-west edge pointing west, its label one clearance band beyond the
% tip, anchored on the paper-east side (the compass involution).
\cs_new_protected:Npn \tenkz_leg_stub:nnNn #1#2#3#4
  {
    % the offsets ride as ++/shift vectors, never inside a calc ($...$):
    % the calc factor parser does not expand a macro-valued coordinate
    \draw[physical~leg] (\l__tenkz_pa_tl.\tenkz_compass:n{#1}) --
      ++\tenkz_leg_vec:nn{#4}{\tenkz@dim{\tenkz@r@physleg}};
    \tl_if_empty:NF #3
      {
        \node[label, anchor=\tenkz_compass:n{#2},
              shift={\tenkz_leg_vec:nn{#4}{ \tenkz@dim{\tenkz@r@physleg}
                       + \tenkz@dim{\tenkz@r@labelclear} }}]
          at (\l__tenkz_pa_tl.\tenkz_compass:n{#1})
          { \tl_use:N #3 };
      }
  }

% -- external labels of bead tensors -----------------------------------
% pos=auto: pick the first free compass direction, in reading order
% south, north, then the diagonals; a direction is occupied when a bond,
% leg, or trace uses it.  `label pos` overrides; `label shift` nudges.
\cs_new_protected:Npn \tenkz_draw_dot_labels:
  {
    \int_step_inline:nn { \l__tenkz_nrows_int }
      {
        \int_step_inline:nnn {1} { \l__tenkz_ncols_int }
          { \tenkz_dot_label:nn {##1}{####1} }
      }
  }

\cs_new_protected:Npn \tenkz_dot_label:nn #1#2
  {
    \prop_get:NnNT \l__tenkz_kind_prop {#1-#2} \l_tmpa_tl
      {
        \bool_lazy_or:nnT
          { \str_if_eq_p:Vn \l_tmpa_tl {tn} }
          { \str_if_eq_p:Vn \l_tmpa_tl {tnc} }
          {
            \tenkz_get:NnN \l__tenkz_opts_prop {#1-#2} \l_tmpb_tl
            \exp_args:NV \tenkz_cell_opts:n \l_tmpb_tl
            \bool_if:nT { \tenkz_label_outside_p: }
              {
                \tenkz_get:NnN \l__tenkz_label_prop {#1-#2} \l_tmpb_tl
                \tl_if_blank:VF \l_tmpb_tl
                  {
                    \tenkz_get:NnN \l__tenkz_kind_prop {#1-#2} \l_tmpa_tl
                    \str_if_eq:VnT \l_tmpa_tl {tnc}
                      { \tl_set:Ne \l_tmpb_tl { \overline{\exp_not:V \l_tmpb_tl} } }
                    \tenkz_dot_label_place:nnV {#1}{#2} \l_tmpb_tl
                  }
              }
          }
      }
  }

% is the resolved skin's label EXTERNAL?  Beads always place their label
% outside; triangle skins do so under the dot policy only (the box policy
% inscribes the label in the triangle when it fits).
\cs_new:Npn \tenkz_label_outside_p:
  {
    \bool_lazy_or_p:nn
      { \str_if_eq_p:Vn \l__tenkz_cshape_tl {dot} }
      { \bool_lazy_and_p:nn
          { \str_if_eq_p:ee { \tenkz@tensorstyle } {dot} }
          { \bool_lazy_or_p:nn
              { \str_if_eq_p:Vn \l__tenkz_cshape_tl {tril} }
              { \str_if_eq_p:Vn \l__tenkz_cshape_tl {trir} } } }
  }

% the auto-quadrant algorithm: the free compass direction of a bead in row
% #1, into tl #2.  Row-level data decides: south when the underside is free,
% north when only the underside is busy, south west when both sides are.
% Interface contention (the nopair genre): when this row would take south
% while the row below would face it with a north label across an unpaired
% interface, both label bands would land in the ONE inter-layer gap — two
% labels may not share a band (label doctrine).  The upper row yields the
% gap and takes its free north west quadrant instead: the diagonal clears
% the up leg and the bond by construction.
\cs_new_protected:Npn \tenkz_auto_quadrant:nN #1#2
  {
    \bool_if:nTF { \tenkz_opendown_p:n {#1} || \tenkz_pair_p:n {#1} }
      {
        \bool_if:nTF { \tenkz_openup_p:n {#1} ||
                       \tenkz_uppaired_p:n {#1} }
          { \tl_set:Nn #2 {south~west} }
          { \tl_set:Nn #2 {north} }
      }
      {
        \bool_if:nTF
          { \int_compare_p:nNn {#1} < { \l__tenkz_nrows_int } &&
            ( \tenkz_opendown_p:n {#1+1} || \tenkz_pair_p:n {#1+1} ) &&
            ! ( \tenkz_openup_p:n {#1+1} || \tenkz_uppaired_p:n {#1+1} ) }
          { \tl_set:Nn #2 {north~west} }
          { \tl_set:Nn #2 {south} }
      }
  }

\tl_new:N \l__tenkz_dot_label_tl
\cs_new_protected:Npn \tenkz_dot_label_place:nnn #1#2#3
  {
    % choose direction.  The auto quadrant reasons in the LOGICAL frame
    % (row occupancy is logical data), so its answer passes through the
    % compass map before the paper-frame placement table below; a
    % user-given `label pos` is already a PAPER word (the reader names
    % what they see) and is used as is.
    \tl_if_empty:NTF \l__tenkz_clpos_tl
      {
        \tenkz_auto_quadrant:nN {#1} \l_tmpa_tl
        \tl_set:Ne \l_tmpa_tl { \tenkz_compass:V \l_tmpa_tl }
      }
      { \tl_set_eq:NN \l_tmpa_tl \l__tenkz_clpos_tl }
    \tenkz_get:NnN \l__tenkz_name_prop {#1-#2} \l__tenkz_pa_tl
    \tl_if_empty:NT \l__tenkz_clshift_tl
      { \tl_set:Nn \l__tenkz_clshift_tl {0pt,0pt} }
    \tl_set:Nn \l__tenkz_dot_label_tl {#3}
    % compass table: node point, opposite label anchor, and the outward
    % (dx,dy); a cardinal nudges one clearance band along its axis, a
    % diagonal 0.7 of a band on each axis
    \str_case:Vn \l_tmpa_tl
      {
        {south}      { \tenkz_dot_label_at:nnnn {south}      {north}
                         {0}{-\tenkz@dim{\tenkz@r@labelclear}} }
        {north}      { \tenkz_dot_label_at:nnnn {north}      {south}
                         {0}{\tenkz@dim{\tenkz@r@labelclear}} }
        {south~west} { \tenkz_dot_label_at:nnnn {south~west} {north~east}
                         {-0.7\tenkz@dim{\tenkz@r@labelclear}}
                         {-0.7\tenkz@dim{\tenkz@r@labelclear}} }
        {south~east} { \tenkz_dot_label_at:nnnn {south~east} {north~west}
                         {0.7\tenkz@dim{\tenkz@r@labelclear}}
                         {-0.7\tenkz@dim{\tenkz@r@labelclear}} }
        {north~west} { \tenkz_dot_label_at:nnnn {north~west} {south~east}
                         {-0.7\tenkz@dim{\tenkz@r@labelclear}}
                         {0.7\tenkz@dim{\tenkz@r@labelclear}} }
        {north~east} { \tenkz_dot_label_at:nnnn {north~east} {south~west}
                         {0.7\tenkz@dim{\tenkz@r@labelclear}}
                         {0.7\tenkz@dim{\tenkz@r@labelclear}} }
      }
  }
\cs_generate_variant:Nn \tenkz_dot_label_place:nnn { nnV }
% place a bead's external label: node compass point #1, opposite label anchor
% #2, outward offset (#3,#4); the text is \l__tenkz_dot_label_tl
\cs_new_protected:Npn \tenkz_dot_label_at:nnnn #1#2#3#4
  {
    \node[label, anchor=#2, shift={(\l__tenkz_clshift_tl)}] at
      ($(\l__tenkz_pa_tl.#1)+(#3,#4)$)
      { $\tenkz@labelsize \tl_use:N \l__tenkz_dot_label_tl$ };
  }

% -- decorations: spans, cuts, bond labels ------------------------------
\cs_new_protected:Npn \tenkz_draw_decorations:
  {
    \seq_map_inline:Nn \l__tenkz_spans_seq
      { \tenkz_draw_span:nnnnn ##1 }
    \seq_map_inline:Nn \l__tenkz_cuts_seq
      { \tenkz_draw_cut:nnn ##1 }
    \seq_map_inline:Nn \l__tenkz_bondlabels_seq
      { \tenkz_draw_bondlabel:nnn ##1 }
  }

% do columns #2..#3 of row #1 place any ACTUAL external dot label on side
% #4 (south|north)?  True only for a dot-skinned tn/tnc cell with a
% non-blank label whose direction (`label pos` override, else the auto
% quadrant) has a #4 component.  Box-skinned cells inscribe their labels,
% so a box-mode span reports false and the brace clearance tightens: no
% phantom band.  The scan is span-local so a brace over unlabeled cells is
% not pushed out by labels elsewhere in the row.
\cs_new_protected:Npn \tenkz_span_dot_labels:nnnnN #1#2#3#4#5
  {
    \bool_set_false:N #5
    \int_step_inline:nnn {#2} {#3}
      {
        \tenkz_get:NeN \l__tenkz_kind_prop {#1-##1} \l_tmpa_tl
        \bool_lazy_or:nnT
          { \str_if_eq_p:Vn \l_tmpa_tl {tn} }
          { \str_if_eq_p:Vn \l_tmpa_tl {tnc} }
          {
            \tenkz_get:NeN \l__tenkz_label_prop {#1-##1} \l_tmpb_tl
            \tl_if_blank:VF \l_tmpb_tl
              {
                \tenkz_get:NeN \l__tenkz_opts_prop {#1-##1} \l_tmpb_tl
                \exp_args:NV \tenkz_cell_opts:n \l_tmpb_tl
                \bool_if:nT { \tenkz_label_outside_p: }
                  {
                    % the side test #4 is a LOGICAL word; a user-given
                    % `label pos` is a PAPER word, so the compass
                    % involution converts it back before comparing
                    \tl_if_empty:NTF \l__tenkz_clpos_tl
                      { \tenkz_auto_quadrant:nN {#1} \l_tmpa_tl }
                      { \tl_set:Ne \l_tmpa_tl
                          { \tenkz_compass:V \l__tenkz_clpos_tl } }
                    \str_if_in:NnT \l_tmpa_tl {#4}
                      { \bool_set_true:N #5 }
                  }
              }
          }
      }
  }

\cs_new_protected:Npn \tenkz_draw_span:nnnnn #1#2#3#4#5
  {
    \dim_set:Nn \l__tenkz_x_dim
      { \tenkz_colx:n {#2} - \tenkz@dim{0.22} }
    \dim_set:Nn \l__tenkz_tmp_dim
      { \tenkz_colx:n { #2 + #3 - 1 } + \tenkz@dim{0.22} }
    \str_if_eq:nnTF {#4} {brace~below}
      {
        % clear open down legs, or the south dot-label band when the row
        % actually places external dot labels below (dot skins only, so a
        % box-mode brace sits closer)
        \tenkz_span_dot_labels:nnnnN {#1} {#2} { #2 + #3 - 1 } {south} \l_tmpa_bool
        \dim_set:Nn \l__tenkz_y_dim
          { \tenkz_rowy:n {#1} - \tenkz@dim{0.20}
            - \bool_if:nTF { \tenkz_opendown_p:n {#1} }
                { \tenkz@dim{\tenkz@r@physleg} }
                { \bool_if:NTF \l_tmpa_bool
                    { \tenkz@dim{\tenkz@r@dotlabelband} } { 0pt } } }
        % brace bulges downward: trace right-to-left, label below
        \tenkz_span_brace:NNnnn \l__tenkz_tmp_dim \l__tenkz_x_dim {north} {-1} {#5}
      }
      {
        % clear open up legs, or the north dot-label band when the row
        % actually places external dot labels above
        \tenkz_span_dot_labels:nnnnN {#1} {#2} { #2 + #3 - 1 } {north} \l_tmpa_bool
        \dim_set:Nn \l__tenkz_y_dim
          { \tenkz_rowy:n {#1} + \tenkz@dim{0.20}
            + \bool_if:nTF { \tenkz_openup_p:n {#1} }
                { \tenkz@dim{\tenkz@r@physleg} }
                { \bool_if:NTF \l_tmpa_bool
                    { \tenkz@dim{\tenkz@r@dotlabelband} } { 0pt } } }
        % brace bulges upward: trace left-to-right, label above
        \tenkz_span_brace:NNnnn \l__tenkz_x_dim \l__tenkz_tmp_dim {south} {1} {#5}
      }
  }

% the shared brace+label of a span, given the two endpoints (their order sets
% which way the brace bulges), the label anchor, the label-offset sign, and
% the label text.  The label sits centred on the span, one text height plus
% one clearance band beyond the brace.
\cs_new_protected:Npn \tenkz_span_brace:NNnnn #1#2#3#4#5
  {
    \draw[brace]
      (\dim_use:N #1, \dim_use:N \l__tenkz_y_dim) --
      (\dim_use:N #2, \dim_use:N \l__tenkz_y_dim);
    \node[label, anchor=#3] at
      (\dim_eval:n{ (\l__tenkz_x_dim + \l__tenkz_tmp_dim)/2 },
       \dim_eval:n{ \l__tenkz_y_dim
                    + ( 3.4pt + \tenkz@dim{\tenkz@r@labelclear} ) * #4 })
      { $\tenkz@labelsize #5$ };
  }

\cs_new_protected:Npn \tenkz_draw_cut:nnn #1#2#3
  {
    \dim_set:Nn \l__tenkz_x_dim
      { \tenkz_colx:n {#2} + 0.5\tenkz@pitch }
    \draw[cut]
      (\dim_use:N \l__tenkz_x_dim,
       \dim_eval:n{ \tenkz_rowy:n {1} + \tenkz@dim{0.55} }) --
      (\dim_use:N \l__tenkz_x_dim,
       \dim_eval:n{ \tenkz_rowy:n {\l__tenkz_nrows_int} - \tenkz@dim{0.55} });
    \tl_if_blank:nF {#3}
      {
        \node[label, anchor=south] at
          (\dim_use:N \l__tenkz_x_dim,
           \dim_eval:n{ \tenkz_rowy:n {1} + \tenkz@dim{0.55}
                        + \tenkz@dim{\tenkz@r@labelclear} })
          { $\tenkz@labelsize #3$ };
      }
  }

\cs_new_protected:Npn \tenkz_draw_bondlabel:nnn #1#2#3
  {
    \node[label, anchor=south] at
      (\dim_eval:n{ ( \tenkz_colx:n {#2} + \tenkz_colx:n {#3} ) / 2 },
       \dim_eval:n{ \tenkz_rowy:n {1} + \tenkz@dim{\tenkz@r@labelclear} })
      { #1 };
  }

% -- open boundaries ----------------------------------------------------
% An open virtual index is drawn as a protruding stub (\tenkz@r@stub): the
% ink states that the object is matrix-valued on that side.  Row-level
% resolution: rows= mod `open`/`none` overrides the picture policy.
\int_new:N \l__tenkz_first_int
\int_new:N \l__tenkz_last_int
\bool_new:N \l__tenkz_bopen_bool
\bool_new:N \l__tenkz_wopen_bool
\bool_new:N \l__tenkz_eopen_bool
\bool_new:N \l__tenkz_wdone_bool
\bool_new:N \l__tenkz_edone_bool

\cs_new_protected:Npn \tenkz_row_extent:n #1
  {
    \int_zero:N \l__tenkz_first_int
    \int_zero:N \l__tenkz_last_int
    \int_step_inline:nnn {1} { \l__tenkz_ncols_int }
      {
        \prop_if_in:NeT \l__tenkz_owner_prop {\int_eval:n{#1}-##1}
          {
            \int_compare:nNnT { \l__tenkz_first_int } = {0}
              { \int_set:Nn \l__tenkz_first_int {##1} }
            \int_set:Nn \l__tenkz_last_int {##1}
          }
      }
  }

\cs_new_protected:Npn \tenkz_row_open:nN #1#2
  {
    \bool_set_false:N #2
    \tl_set:Ne \l_tmpb_tl { \tenkz_rowmod:n {#1} }
    \str_if_in:NnTF \l_tmpb_tl {open}
      { \bool_set_true:N #2 }
      {
        \str_if_in:NnF \l_tmpb_tl {none}
          { \str_if_eq:VnT \l__tenkz_boundary_tl {open}
              { \bool_set_true:N #2 } }
      }
  }

% per-row, per-side openness: is row #1 open on side #2 (west|east)?
% Resolution order: the rows= suffix wins — the side-scoped `west open`
% / `west none` / `east open` / `east none` first, then the both-sides
% `open` / `none` — then the side policy (west=/east=), then the
% picture policy (boundary=).  Prepared and discarded indices are SIDES
% of the word, not rows -- a circuit fragment seals its west (input)
% indices and keeps its east (output) indices open (hard11) -- so a row
% whose west resolves to `none` draws no west stub yet keeps its east
% stub, and the boundary signature counts only the stubs actually
% drawn.  The row-scope words live in \l__tenkz_rowside_prop ("side-r"
% -> open|none), built once per picture by \tenkz_record_rowsides: from
% the exact colon-split mod segments (a substring test would let `west
% open` open BOTH sides through its `open` letters).
\cs_new_protected:Npn \tenkz_row_open_side:nnN #1#2#3
  {
    \bool_set_false:N #3
    \prop_get:NeNTF \l__tenkz_rowside_prop { #2 - \int_eval:n {#1} }
      \l_tmpb_tl
      { \str_if_eq:VnT \l_tmpb_tl {open} { \bool_set_true:N #3 } }
      {
        \str_if_eq:nnTF {#2} {west}
          { \tl_set_eq:NN \l_tmpb_tl \l__tenkz_westpol_tl }
          { \tl_set_eq:NN \l_tmpb_tl \l__tenkz_eastpol_tl }
        \tl_if_empty:NTF \l_tmpb_tl
          { \str_if_eq:VnT \l__tenkz_boundary_tl {open}
              { \bool_set_true:N #3 } }
          { \str_if_eq:VnT \l_tmpb_tl {open}
              { \bool_set_true:N #3 } }
      }
  }

% build the row-scope side-policy table: one pass over every row's mod
% segments (split at the catcode-12 colons of the str-normalized mods),
% spaces stripped, exact words only — `open`/`none` write both sides,
% `west open` etc. one side (`rows={op:west none}`: the op row seals its
% west end and keeps its east stub).  Later segments override earlier
% ones, matching the option-list read.
\prop_new:N \l__tenkz_rowside_prop
\seq_new:N \l__tenkz_mods_seq
\cs_new_protected:Npn \tenkz_record_rowsides:
  {
    \prop_clear:N \l__tenkz_rowside_prop
    \int_step_inline:nn { \l__tenkz_nrows_int }
      {
        \tl_set:Ne \l_tmpb_tl { \tenkz_rowmod:n {##1} }
        \use:e
          {
            \seq_set_split:Nnn \exp_not:N \l__tenkz_mods_seq
              { \c_colon_str } { \exp_not:V \l_tmpb_tl }
          }
        \seq_map_inline:Nn \l__tenkz_mods_seq
          { \tenkz_rowside_mod:nn {##1} {####1} }
      }
  }
\cs_new_protected:Npn \tenkz_rowside_mod:nn #1#2
  {
    \tl_set:Nn \l_tmpb_tl {#2}
    \tl_remove_all:Nn \l_tmpb_tl { ~ }
    \str_case:VnF \l_tmpb_tl
      {
        {open}     { \prop_put:Nen \l__tenkz_rowside_prop {west-#1} {open}
                     \prop_put:Nen \l__tenkz_rowside_prop {east-#1} {open} }
        {none}     { \prop_put:Nen \l__tenkz_rowside_prop {west-#1} {none}
                     \prop_put:Nen \l__tenkz_rowside_prop {east-#1} {none} }
        {westopen} { \prop_put:Nen \l__tenkz_rowside_prop {west-#1} {open} }
        {westnone} { \prop_put:Nen \l__tenkz_rowside_prop {west-#1} {none} }
        {eastopen} { \prop_put:Nen \l__tenkz_rowside_prop {east-#1} {open} }
        {eastnone} { \prop_put:Nen \l__tenkz_rowside_prop {east-#1} {none} }
      }
      { }
  }

% is the cell at (row, col) part of a fusion bar?  A bar consumes the
% boundary on its side: its combined leg IS the open index of the fused
% object, so no stub is drawn behind it.
\cs_new_protected:Npn \tenkz_cell_is_fuse:nnN #1#2#3
  {
    \bool_set_false:N #3
    \tenkz_get:NeN \l__tenkz_kind_prop {\int_eval:n{#1}-\int_eval:n{#2}} \l_tmpa_tl
    \bool_lazy_or:nnT
      { \str_if_eq_p:Vn \l_tmpa_tl {fuse} }
      { \str_if_eq_p:Vn \l_tmpa_tl {fusec} }
      { \bool_set_true:N #3 }
  }

\cs_new_protected:Npn \tenkz_draw_boundaries:
  {
    \int_zero:N \l__tenkz_bvw_int
    \int_zero:N \l__tenkz_bve_int
    \bool_set_false:N \l__tenkz_wdone_bool
    \bool_set_false:N \l__tenkz_edone_bool
    \int_step_inline:nn { \l__tenkz_nrows_int }
      {
        % each side resolves its own openness (rows= mod, then west=/
        % east=, then boundary=): a west=none row keeps its east stub
        \tenkz_row_open_side:nnN {##1}{west} \l__tenkz_wopen_bool
        \tenkz_row_open_side:nnN {##1}{east} \l__tenkz_eopen_bool
        \bool_lazy_or:nnT { \l__tenkz_wopen_bool } { \l__tenkz_eopen_bool }
          {
            \tenkz_row_extent:n {##1}
            \int_compare:nNnF { \l__tenkz_first_int } = {0}
              {
                \tenkz_row_wirestyle:nN {##1} \l__tenkz_cell_tl
                % west stub, unless the side is closed, a fusion bar owns
                % this end, or a cup closes it (a cupped end is
                % contracted, not open)
                \bool_if:NT \l__tenkz_wopen_bool
                  {
                \tenkz_cell_is_fuse:nnN {##1} {\l__tenkz_first_int}
                  \l_tmpa_bool
                \bool_if:NF \l_tmpa_bool
                  {
                  \prop_if_in:NeF \l__tenkz_cuprow_prop { west - ##1 }
                  {
                    \tenkz_anchor:nnnN {##1}{\int_use:N \l__tenkz_first_int}
                      {west} \l__tenkz_pa_tl
                    \draw[\tl_use:N \l__tenkz_cell_tl]
                      \l__tenkz_pa_tl --
                      ++\tenkz_stub_vec:n{west}
                      coordinate (tzbw\the\tenkz@pictureid);
                    \int_incr:N \l__tenkz_bvw_int
                    % the west boundary label attaches to the first row
                    % that draws a west stub, one clearance band beyond
                    % the tip along the stub's own (frame-mapped) axis
                    \bool_if:NF \l__tenkz_wdone_bool
                      {
                        \bool_set_true:N \l__tenkz_wdone_bool
                        \tl_if_empty:NF \l__tenkz_wlabel_tl
                          {
                            \node[label, anchor=\tenkz_compass:n{east},
                                  shift={\tenkz_side_vec:nn{west}
                                    {\tenkz@dim{\tenkz@r@labelclear}}}]
                              at (tzbw\the\tenkz@pictureid)
                              { \tl_use:N \l__tenkz_wlabel_tl };
                          }
                      }
                  }
                  }
                  }
                % east stub, same rules
                \bool_if:NT \l__tenkz_eopen_bool
                  {
                \tenkz_cell_is_fuse:nnN {##1} {\l__tenkz_last_int}
                  \l_tmpa_bool
                \bool_if:NF \l_tmpa_bool
                  {
                  \prop_if_in:NeF \l__tenkz_cuprow_prop { east - ##1 }
                  {
                    \tenkz_anchor:nnnN {##1}{\int_use:N \l__tenkz_last_int}
                      {east} \l__tenkz_pa_tl
                    \draw[\tl_use:N \l__tenkz_cell_tl]
                      \l__tenkz_pa_tl --
                      ++\tenkz_stub_vec:n{east}
                      coordinate (tzbe\the\tenkz@pictureid);
                    \int_incr:N \l__tenkz_bve_int
                    \bool_if:NF \l__tenkz_edone_bool
                      {
                        \bool_set_true:N \l__tenkz_edone_bool
                        \tl_if_empty:NF \l__tenkz_elabel_tl
                          {
                            \node[label, anchor=\tenkz_compass:n{west},
                                  shift={\tenkz_side_vec:nn{east}
                                    {\tenkz@dim{\tenkz@r@labelclear}}}]
                              at (tzbe\the\tenkz@pictureid)
                              { \tl_use:N \l__tenkz_elabel_tl };
                          }
                      }
                  }
                  }
                  }
              }
          }
      }
    % a \tnfuse combined stub is ONE open virtual index on its side: the
    % fusion map V : C^{D_1} x C^{D_2} -> C^{D_12} leaves exactly the fused
    % index open, and the combined leg IS that index.  Without this count a
    % fused equation reads unbalanced -- hard06's O = X^{-1}[A (x) conj A]X
    % needs virtual-west = virtual-east = 1 on BOTH sides.
    \int_step_inline:nn { \l__tenkz_nrows_int }
      {
        \int_step_inline:nnn {1} { \l__tenkz_ncols_int }
          {
            \tenkz_get:NeN \l__tenkz_kind_prop {##1-####1} \l__tenkz_pkind_tl
            \str_if_eq:VnT \l__tenkz_pkind_tl {fuse}
              {
                \tenkz_get:NeN \l__tenkz_opts_prop {##1-####1} \l_tmpb_tl
                \exp_args:NV \tenkz_cell_opts:n \l_tmpb_tl
                \str_if_eq:VnTF \l__tenkz_fcombined_tl {east}
                  { \int_incr:N \l__tenkz_bve_int }
                  { \int_incr:N \l__tenkz_bvw_int }
              }
          }
      }
  }

% -- inter-layer cups -----------------------------------------------------
% A cup closes two DIFFERENT rows into each other on one side — the
% ket-to-bra virtual tie of sandwich objects.  Contrast periodic: a trace
% closes one row back to ITSELF; the two closures have different topology
% and the event stream records them as different events (trace vs cup).
%
% Pair selection.  Eligible rows on a side are the rows that would draw a
% stub there: open under the row/picture boundary policy, non-empty, and
% not ending in a fusion bar (the bar consumes its side's boundary).
% Eligible rows are cupped in ADJACENT PAIRS, top-down (first with second,
% third with fourth, ...); an odd row out keeps its stub.  A cupped row
% draws no stub and contributes no open virtual index to the signature.
\cs_new_protected:Npn \tenkz_compute_cups:
  {
    \seq_clear:N \l__tenkz_cupw_seq
    \seq_clear:N \l__tenkz_cupe_seq
    \prop_clear:N \l__tenkz_cuprow_prop
    \bool_if:NT \l__tenkz_closew_bool
      { \tenkz_compute_cups_side:nnN {west}{first} \l__tenkz_cupw_seq }
    \bool_if:NT \l__tenkz_closee_bool
      { \tenkz_compute_cups_side:nnN {east}{last} \l__tenkz_cupe_seq }
  }
% side #1 (west|east), row-extent end #2 (first|last), pair seq #3
\cs_new_protected:Npn \tenkz_compute_cups_side:nnN #1#2#3
  {
    \seq_clear:N \l_tmpa_seq
    \int_step_inline:nn { \l__tenkz_nrows_int }
      {
        % cup eligibility respects the per-side policy: a side sealed by
        % west=none/east=none has no open ends to cup there
        \tenkz_row_open_side:nnN {##1}{#1} \l__tenkz_bopen_bool
        \bool_if:NT \l__tenkz_bopen_bool
          {
            \tenkz_row_extent:n {##1}
            \int_compare:nNnF { \l__tenkz_first_int } = {0}
              {
                \tenkz_cell_is_fuse:nnN {##1}
                  { \int_use:c { l__tenkz_#2_int } } \l_tmpa_bool
                \bool_if:NF \l_tmpa_bool
                  { \seq_put_right:Nn \l_tmpa_seq {##1} }
              }
          }
      }
    \bool_while_do:nn
      { \int_compare_p:nNn { \seq_count:N \l_tmpa_seq } > {1} }
      {
        \seq_pop_left:NN \l_tmpa_seq \l_tmpa_tl
        \seq_pop_left:NN \l_tmpa_seq \l_tmpb_tl
        \seq_put_right:Ne #3 { { \l_tmpa_tl } { \l_tmpb_tl } }
        \prop_put:Nen \l__tenkz_cuprow_prop { #1 - \l_tmpa_tl } {}
        \prop_put:Nen \l__tenkz_cuprow_prop { #1 - \l_tmpb_tl } {}
      }
  }

% Geometry.  One cup = two quarter-arcs (a rounded 180-degree turn) from
% the upper row's end to the lower row's end, drawn in the pair's wire
% style (the upper row's — a sandwich pair shares one style).  Horizontal
% reach = \tenkz@r@stub: the bend protrudes exactly as far as the stubs it
% replaces, so open and cup-closed pictures share one silhouette; being
% purely outboard of the last column it also clears every leg-tip and dot
% label band by construction (labels live at column abscissae).
\dim_new:N \l__tenkz_cupx_dim   % common abscissa where the turn starts
\dim_new:N \l__tenkz_cupry_dim  % vertical radius: half the pair's row gap
\tl_new:N \l__tenkz_cupmat_tl   % this cup's on-bend matrix (may be empty)

\cs_new_protected:Npn \tenkz_draw_cups:
  {
    \seq_map_inline:Nn \l__tenkz_cupw_seq
      { \tenkz_draw_cup:nnn {west} ##1 }
    \seq_map_inline:Nn \l__tenkz_cupe_seq
      { \tenkz_draw_cup:nnn {east} ##1 }
  }

% the cup end of row #2 on side #1: the anchor coordinate (into tl #3) and
% its abscissa (into dim #4).  Cup ends are always glyph nodes — a fuse-
% consumed end is never cup-eligible — so the abscissa comes straight from
% the pgf shape anchor.
\cs_new_protected:Npn \tenkz_cup_anchor:nnNN #1#2#3#4
  {
    \tenkz_row_extent:n {#2}
    \str_if_eq:nnTF {#1}{west}
      { \int_set:Nn \l_tmpb_int { \l__tenkz_first_int } }
      { \int_set:Nn \l_tmpb_int { \l__tenkz_last_int } }
    \tenkz_anchor:nnnN {#2}{ \int_use:N \l_tmpb_int }{#1} #3
    \prop_get:NeNF \l__tenkz_owner_prop
      {#2-\int_use:N \l_tmpb_int} \l_tmpa_tl
      { \tl_set:Ne \l_tmpa_tl { \int_use:N \l_tmpb_int } }
    \tenkz_get:NeN \l__tenkz_name_prop {#2-\l_tmpa_tl} \l_tmpb_tl
    \pgfextractx #4 { \exp_args:Ne \pgfpointanchor { \l_tmpb_tl } {#1} }
  }

% draw one cup: side #1 (west|east), upper row #2, lower row #3
\cs_new_protected:Npn \tenkz_draw_cup:nnn #1#2#3
  {
    \tenkz_row_wirestyle:nN {#2} \l__tenkz_cell_tl
    \str_if_eq:nnTF {#1}{east}
      { \tl_set_eq:NN \l__tenkz_cupmat_tl \l__tenkz_closeemat_tl }
      { \tl_set_eq:NN \l__tenkz_cupmat_tl \l__tenkz_closewmat_tl }
    \tenkz_cup_anchor:nnNN {#1}{#2} \l__tenkz_pa_tl \l__tenkz_x_dim
    \tenkz_cup_anchor:nnNN {#1}{#3} \l__tenkz_pb_tl \l__tenkz_tmp_dim
    % common abscissa: the OUTERMOST of the two ends; a shorter row gets a
    % horizontal lead-in so the turn itself stays one smooth half-ellipse
    \str_if_eq:nnTF {#1}{east}
      { \dim_set:Nn \l__tenkz_cupx_dim
          { \dim_max:nn { \l__tenkz_x_dim } { \l__tenkz_tmp_dim } } }
      { \dim_set:Nn \l__tenkz_cupx_dim
          { \dim_min:nn { \l__tenkz_x_dim } { \l__tenkz_tmp_dim } } }
    \dim_set:Nn \l__tenkz_cupry_dim
      { ( \tenkz_rowy:n {#2} - \tenkz_rowy:n {#3} ) / 2 }
    \dim_set:Nn \l__tenkz_y_dim
      { ( \tenkz_rowy:n {#2} + \tenkz_rowy:n {#3} ) / 2 }
    \str_if_eq:nnTF {#1}{east}
      {
        \draw[\tl_use:N \l__tenkz_cell_tl]
          \l__tenkz_pa_tl --
          (\dim_use:N \l__tenkz_cupx_dim, \tenkz_rowy:n {#2})
          arc[start~angle=90, end~angle=0,
              x~radius=\tenkz@dim{\tenkz@r@stub},
              y~radius=\dim_use:N \l__tenkz_cupry_dim]
          arc[start~angle=0, end~angle=-90,
              x~radius=\tenkz@dim{\tenkz@r@stub},
              y~radius=\dim_use:N \l__tenkz_cupry_dim]
          -- \l__tenkz_pb_tl;
        \tenkz_cup_matrix:nnn {#1}{#2}{ 1 }
      }
      {
        \draw[\tl_use:N \l__tenkz_cell_tl]
          \l__tenkz_pa_tl --
          (\dim_use:N \l__tenkz_cupx_dim, \tenkz_rowy:n {#2})
          arc[start~angle=90, end~angle=180,
              x~radius=\tenkz@dim{\tenkz@r@stub},
              y~radius=\dim_use:N \l__tenkz_cupry_dim]
          arc[start~angle=180, end~angle=270,
              x~radius=\tenkz@dim{\tenkz@r@stub},
              y~radius=\dim_use:N \l__tenkz_cupry_dim]
          -- \l__tenkz_pb_tl;
        \tenkz_cup_matrix:nnn {#1}{#2}{ -1 }
      }
    \tenkz@event{cup|picture=\the\tenkz@pictureid|side=#1|top=#2|bottom=#3|matrix=\tl_if_empty:NTF \l__tenkz_cupmat_tl {0}{1}}
  }

% the on-bend fixed-point matrix: a filled dot ON the bend apex, its label
% OUTSIDE the bend, one clearance band beyond the dot on the side the bend
% points to — the rho^L / rho^R capping convention (RMP sec2fig8, fig9-1).
% #1 side, #2 upper row (they name the node), #3 outward sign (1|-1)
\cs_new_protected:Npn \tenkz_cup_matrix:nnn #1#2#3
  {
    \tl_if_blank:VF \l__tenkz_cupmat_tl
      {
        \node[tensor]
          (tzcup\the\tenkz@pictureid-#1-#2)
          at (\dim_eval:n{ \l__tenkz_cupx_dim
                           + \tenkz@dim{\tenkz@r@stub} * (#3) },
              \dim_use:N \l__tenkz_y_dim) {};
        \str_if_eq:nnTF {#1}{east}
          {
            \node[label, anchor=west] at
              ($(tzcup\the\tenkz@pictureid-#1-#2.east)
                +(\tenkz@dim{\tenkz@r@labelclear},0)$)
              { \tl_use:N \l__tenkz_cupmat_tl };
          }
          {
            \node[label, anchor=east] at
              ($(tzcup\the\tenkz@pictureid-#1-#2.west)
                +(-\tenkz@dim{\tenkz@r@labelclear},0)$)
              { \tl_use:N \l__tenkz_cupmat_tl };
          }
      }
  }

% The boundary signature counts the picture's open indices.  Two sides of
% a diagrammatic equation are equal objects only if their signatures
% match; the audit layer reads these events to lint that.
% the number of open legs a leg-candidate cell contributes: 1, or the
% number of `legs at` columns (each named column is its own physical index)
\cs_generate_variant:Nn \clist_count:n { V }
\cs_new:Npn \tenkz_leg_count:
  {
    \tl_if_empty:NTF \l__tenkz_clegsat_tl {1}
      { \clist_count:V \l__tenkz_clegsat_tl }
  }

\cs_new_protected:Npn \tenkz_emit_signature:
  {
    \int_zero:N \l__tenkz_bpu_int
    \int_zero:N \l__tenkz_bpd_int
    % trace=physical closes each traced site's up+down pair per site, so
    % those legs leave the signature (single-row form only; the deferred
    % multi-row form keeps its counts until it draws)
    \bool_lazy_and:nnTF { \l__tenkz_ptrace_bool }
      { \int_compare_p:nNn { \l__tenkz_nrows_int } = {1} }
      { \bool_set_true:N  \l__tenkz_ptraceon_bool }
      { \bool_set_false:N \l__tenkz_ptraceon_bool }
    \int_step_inline:nn { \l__tenkz_nrows_int }
      {
        \int_step_inline:nnn {1} { \l__tenkz_ncols_int }
          {
            \bool_if:NF \l__tenkz_ptraceon_bool
              {
                % a traced column contributes NO open physical index: the
                % wrap loop closes ket against bra (Tr_s), so both the ket
                % cell and the bra cell below it leave the signature
                \bool_lazy_and:nnT
                  { \tenkz_openup_p:n {##1} }
                  { ! \tenkz_traced_p:nn {##1}{####1} }
                  { \tenkz_leg_candidate:nnT {##1}{####1}
                      { \int_add:Nn \l__tenkz_bpu_int { \tenkz_leg_count: } } }
                \bool_lazy_and:nnT
                  { \tenkz_opendown_p:n {##1} }
                  { ! \tenkz_traced_above_p:nn {##1}{####1} }
                  { \tenkz_leg_candidate:nnT {##1}{####1}
                      { \int_add:Nn \l__tenkz_bpd_int { \tenkz_leg_count: } } }
              }
            % a \tnskip hole at a paired interface leaves the live partner's
            % leg OPEN: that leg is an open physical index and must count
            \int_compare:nNnT {##1} < { \l__tenkz_nrows_int }
              {
                \bool_if:nT { \tenkz_pair_p:n {##1} }
                  {
                    \tenkz_get:NeN \l__tenkz_kind_prop
                      { \int_eval:n{##1+1} - ####1 } \l__tenkz_pkind_tl
                    \str_if_eq:VnT \l__tenkz_pkind_tl {skip}
                      { \tenkz_leg_candidate:nnT {##1}{####1}
                          { \int_add:Nn \l__tenkz_bpd_int
                              { \tenkz_leg_count: } } }
                    \tenkz_get:NeN \l__tenkz_kind_prop
                      {##1-####1} \l__tenkz_pkind_tl
                    \str_if_eq:VnT \l__tenkz_pkind_tl {skip}
                      { \tenkz_leg_candidate:nnT
                          { \int_eval:n{##1+1} }{####1}
                          { \int_add:Nn \l__tenkz_bpu_int
                              { \tenkz_leg_count: } } }
                    % an OPENED interface segment (open={(i,c)}) keeps
                    % both faces' stubs: two open physical indices
                    \prop_if_in:NeT \l__tenkz_opencell_prop {##1-####1}
                      {
                        \tenkz_leg_candidate:nnT {##1}{####1}
                          { \int_add:Nn \l__tenkz_bpd_int
                              { \tenkz_leg_count: } }
                        \tenkz_leg_candidate:nnT
                          { \int_eval:n{##1+1} }{####1}
                          { \int_add:Nn \l__tenkz_bpu_int
                              { \tenkz_leg_count: } }
                      }
                  }
              }
          }
      }
    \tenkz@event{boundary|picture=\the\tenkz@pictureid|virtual-west=\int_use:N\l__tenkz_bvw_int|virtual-east=\int_use:N\l__tenkz_bve_int|physical-up=\int_use:N\l__tenkz_bpu_int|physical-down=\int_use:N\l__tenkz_bpd_int}
  }

% -- physical trace (trace=physical) --------------------------------------
% Per site, Tr_phys(M) = sum_s M^{ss}: the site's up leg turns east at its
% tip, swings around the glyph's east side at \tenkz@r@phtracereach beyond
% the glyph edge (crossing the virtual wire -- distinct ink, no junction),
% and comes back into the site's own down-leg tip.  The legs themselves are
% drawn by the leg pass; this pass adds only the connecting loop, in the
% `trace` style (solid, rounded corners -- dotted-free).  The multi-row form
% (outermost up leg to outermost down leg around the east end of the WHOLE
% word) is deferred: it warns and draws nothing.
\msg_new:nnn { tenkz } { phtrace-multirow }
  {
    trace=physical~currently~closes~single-row~pictures~only;~the~
    multi-row~form~(outermost~up~leg~to~outermost~down~leg~around~the~
    east~end~of~the~word)~is~deferred.~No~trace~ink~was~drawn.
  }

\cs_new_protected:Npn \tenkz_draw_phtrace:
  {
    \bool_if:NT \l__tenkz_ptrace_bool
      {
        \int_compare:nNnTF { \l__tenkz_nrows_int } = {1}
          {
            \int_step_inline:nnn {1} { \l__tenkz_ncols_int }
              {
                \bool_lazy_and:nnT
                  { \tenkz_openup_p:n {1} }
                  { \tenkz_opendown_p:n {1} }
                  { \tenkz_leg_candidate:nnT {1}{##1}
                      { \tenkz_phtrace_site:n {##1} } }
              }
          }
          { \msg_warning:nn {tenkz}{phtrace-multirow} }
      }
  }

% one site's trace loop: up-leg tip -> around the east side -> down-leg tip
\cs_new_protected:Npn \tenkz_phtrace_site:n #1
  {
    \tenkz_get:NnN \l__tenkz_name_prop {1-#1} \l__tenkz_pa_tl
    % up-leg tip (edge ordinate + leg), down-leg tip, and the loop abscissa
    \pgfextracty \l__tenkz_tmp_dim
      { \exp_args:Ne \pgfpointanchor { \l__tenkz_pa_tl } {north} }
    \dim_add:Nn \l__tenkz_tmp_dim { \tenkz@dim{\tenkz@r@physleg} }
    \pgfextracty \l__tenkz_y_dim
      { \exp_args:Ne \pgfpointanchor { \l__tenkz_pa_tl } {south} }
    \dim_sub:Nn \l__tenkz_y_dim { \tenkz@dim{\tenkz@r@physleg} }
    \pgfextractx \l__tenkz_x_dim
      { \exp_args:Ne \pgfpointanchor { \l__tenkz_pa_tl } {east} }
    \dim_add:Nn \l__tenkz_x_dim { \tenkz@dim{\tenkz@r@phtracereach} }
    % legs live on the glyph's vertical centre line
    \pgfextractx \l__tenkz_tallx_dim
      { \exp_args:Ne \pgfpointanchor { \l__tenkz_pa_tl } {north} }
    \draw[trace]
      (\dim_use:N \l__tenkz_tallx_dim, \dim_use:N \l__tenkz_tmp_dim) --
      (\dim_use:N \l__tenkz_x_dim,     \dim_use:N \l__tenkz_tmp_dim) --
      (\dim_use:N \l__tenkz_x_dim,     \dim_use:N \l__tenkz_y_dim) --
      (\dim_use:N \l__tenkz_tallx_dim, \dim_use:N \l__tenkz_y_dim);
    \tenkz@event{phtrace|picture=\the\tenkz@pictureid|row=1|col=#1}
  }

% -- pair trace (the trace= physical cell-set) ----------------------------
% Per traced column, the wrap loop IS the partial trace over that site's
% physical space: Tr_s contracts the ket's physical index with the bra's,
% and the ink shows the contraction as ONE wire -- up out of the ket, a
% stadium cap around the column's east side, down past both rows, and into
% the bra's down leg from below -- so the traced column leaves no open
% physical index and the picture depicts rho_kept = Tr_traced |psi><psi|
% (RMP 2011.12127 fig. 2 marginals).  The east straight swings
% \tenkz@r@pairtracereach beyond the leg axis, clear of the neighbouring
% column at the default pitch, crossing the virtual wires as distinct ink
% (no junction); the caps are circular (radius = reach/2), so every join
% is tangent-continuous.
\cs_new_protected:Npn \tenkz_draw_pairtrace:
  {
    \int_step_inline:nn { \l__tenkz_nrows_int }
      {
        \int_step_inline:nnn {1} { \l__tenkz_ncols_int }
          {
            \prop_if_in:NeT \l__tenkz_trcell_prop {##1-####1}
              {
                \tenkz_leg_candidate:nnT {##1}{####1}
                  {
                    \tenkz_leg_candidate:nnT { \int_eval:n{##1+1} }{####1}
                      { \tenkz_pairtrace_site:nn {##1}{####1} }
                  }
              }
          }
      }
  }

% one column's wrap loop: ket up-leg tip -> stadium cap east -> down past
% both rows -> stadium cap west -> bra down-leg tip.  Both legs are drawn
% here as part of the loop (a traced cell's stub/pairing passes stand
% down), so the whole trace is one stroke of the `trace` style.
\cs_new_protected:Npn \tenkz_pairtrace_site:nn #1#2
  {
    \tenkz_get:NnN \l__tenkz_name_prop {#1-#2} \l__tenkz_pa_tl
    \tenkz_get:NeN \l__tenkz_name_prop { \int_eval:n {#1+1} - #2 }
      \l__tenkz_pb_tl
    % ket up-leg tip ordinate (edge + leg) and bra down-leg tip ordinate
    \pgfextracty \l__tenkz_tmp_dim
      { \exp_args:Ne \pgfpointanchor { \l__tenkz_pa_tl } {north} }
    \dim_add:Nn \l__tenkz_tmp_dim { \tenkz@dim{\tenkz@r@physleg} }
    \pgfextracty \l__tenkz_y_dim
      { \exp_args:Ne \pgfpointanchor { \l__tenkz_pb_tl } {south} }
    \dim_sub:Nn \l__tenkz_y_dim { \tenkz@dim{\tenkz@r@physleg} }
    % the leg axis (the glyphs' vertical centre line) and the east straight
    \pgfextractx \l__tenkz_tallx_dim
      { \exp_args:Ne \pgfpointanchor { \l__tenkz_pa_tl } {north} }
    \dim_set:Nn \l__tenkz_x_dim
      { \l__tenkz_tallx_dim + \tenkz@dim{\tenkz@r@pairtracereach} }
    \draw[trace]
      (\l__tenkz_pa_tl.north) --
      (\dim_use:N \l__tenkz_tallx_dim, \dim_use:N \l__tenkz_tmp_dim)
      arc[start~angle=180, end~angle=0,
          radius=\dim_eval:n { \tenkz@dim{\tenkz@r@pairtracereach} / 2 }]
      --
      (\dim_use:N \l__tenkz_x_dim, \dim_use:N \l__tenkz_y_dim)
      arc[start~angle=360, end~angle=180,
          radius=\dim_eval:n { \tenkz@dim{\tenkz@r@pairtracereach} / 2 }]
      -- (\l__tenkz_pb_tl.south);
    \tenkz@event{pairtrace|picture=\the\tenkz@pictureid|row=\int_eval:n{#1}|col=#2}
  }

% -- periodic closure ---------------------------------------------------
\cs_new_protected:Npn \tenkz_draw_trace:
  {
    \int_compare:nNnTF { \l__tenkz_nrows_int } = {1}
      { \tenkz_trace_row:nn {1}{below} }
      {
        \tenkz_trace_row:nn {1}{above}
        \tenkz_trace_row:nn {\l__tenkz_nrows_int}{below}
      }
  }

\cs_new_protected:Npn \tenkz_trace_row:nn #1#2
  {
    % closures inherit (hue contract, clause 4): the racetrack draws in
    % the traced row's resolved wire style over the trace geometry, so
    % an operator-hued row closes with an operator-hued ring
    \tenkz_row_wirestyle:nN {#1} \l__tenkz_cell_tl
    % periodic + dir: the barb rides the racetrack's RETURN wire (the MPO
    % ring convention).  The path is stroked from the east end around to the
    % west end, so its long straight already runs with the circulation of
    % right-directed bonds; `left` reverses the barb, not the path.
    \tenkz_row_dir:nN {#1} \l__tenkz_dir_tl
    \tl_clear:N \l__tenkz_dirmarks_tl
    \tl_if_empty:NF \l__tenkz_dir_tl
      {
        \tl_set:Nn \l__tenkz_dirmarks_tl {dir~marks}
        \str_if_eq:VnT \l__tenkz_dir_tl {left}
          { \tl_put_right:Nn \l__tenkz_dirmarks_tl {~reversed} }
      }
    % find first and last occupied columns of the row
    \int_zero:N \l_tmpa_int  \int_zero:N \l_tmpb_int
    \int_step_inline:nnn {1} { \l__tenkz_ncols_int }
      {
        \prop_if_in:NeT \l__tenkz_owner_prop {\int_eval:n{#1}-##1}
          {
            \int_compare:nNnT { \l_tmpa_int } = {0}
              { \int_set:Nn \l_tmpa_int {##1} }
            \int_set:Nn \l_tmpb_int {##1}
          }
      }
    \int_compare:nNnF { \l_tmpa_int } = {0}
      {
        \tenkz_anchor:nnnN {#1}{\int_use:N\l_tmpb_int}{east} \l__tenkz_pa_tl
        \tenkz_anchor:nnnN {#1}{\int_use:N\l_tmpa_int}{west} \l__tenkz_pb_tl
        \dim_set:Nn \l__tenkz_x_dim
          { \tenkz_colx:n {\l_tmpb_int} + \tenkz@dim{0.45} }
        \dim_set:Nn \l__tenkz_tmp_dim
          { \tenkz_colx:n {\l_tmpa_int} - \tenkz@dim{0.45} }
        % The racetrack must CLEAR any ink on its side of the traced row:
        % open physical legs add their length, and leg-tip/dot labels add
        % the label band (a virtual trace crossing physical ink would draw
        % a false contraction).
        \dim_set:Nn \l__tenkz_y_dim { \tenkz@dim{\tenkz@r@traceclear} }
        \str_if_eq:nnTF {#2} {below}
          {
            \bool_if:nT { \tenkz_opendown_p:n {#1} }
              { \dim_add:Nn \l__tenkz_y_dim { \tenkz@dim{\tenkz@r@physleg} } }
            \tenkz_span_dot_labels:nnnnN {#1} {1} {\l__tenkz_ncols_int}
              {south} \l_tmpa_bool
            \bool_if:NT \l_tmpa_bool
              { \dim_add:Nn \l__tenkz_y_dim
                  { \tenkz@dim{\tenkz@r@dotlabelband} } }
            \dim_set:Nn \l__tenkz_y_dim
              { \tenkz_rowy:n {#1} - \l__tenkz_y_dim }
          }
          {
            \bool_if:nT { \tenkz_openup_p:n {#1} }
              { \dim_add:Nn \l__tenkz_y_dim { \tenkz@dim{\tenkz@r@physleg} } }
            \tenkz_span_dot_labels:nnnnN {#1} {1} {\l__tenkz_ncols_int}
              {north} \l_tmpa_bool
            \bool_if:NT \l_tmpa_bool
              { \dim_add:Nn \l__tenkz_y_dim
                  { \tenkz@dim{\tenkz@r@dotlabelband} } }
            \dim_set:Nn \l__tenkz_y_dim
              { \tenkz_rowy:n {#1} + \l__tenkz_y_dim }
          }
        % the four racetrack corners are logical-frame points; each one
        % passes through the frame map, so a vertical picture gets the
        % side racetrack (the return wire runs down the page's east or
        % west margin) from the same clearance arithmetic
        \tl_set:Ne \l_tmpa_tl { \exp_not:V \l__tenkz_pa_tl }
        \tenkz_map:nn { \l__tenkz_x_dim } { \tenkz_rowy:n {#1} }
        \tl_put_right:Ne \l_tmpa_tl { -- \tenkz_mpt: }
        \tenkz_map:nn { \l__tenkz_x_dim } { \l__tenkz_y_dim }
        \tl_put_right:Ne \l_tmpa_tl { -- \tenkz_mpt: }
        \tenkz_map:nn { \l__tenkz_tmp_dim } { \l__tenkz_y_dim }
        \tl_put_right:Ne \l_tmpa_tl { -- \tenkz_mpt: }
        \tenkz_map:nn { \l__tenkz_tmp_dim } { \tenkz_rowy:n {#1} }
        \tl_put_right:Ne \l_tmpa_tl { -- \tenkz_mpt: }
        \draw[trace, \tl_use:N \l__tenkz_cell_tl,
              \tl_use:N \l__tenkz_dirmarks_tl]
          \l_tmpa_tl -- \l__tenkz_pb_tl;
        % row=... must write a NUMBER: #1 may be the nrows register, and
        % an unexpandable register token would land verbatim in the log
        \tenkz@event{trace|picture=\the\tenkz@pictureid|row=\int_eval:n{#1}|side=#2}
      }
  }

% -- periodic closure, hooks style --------------------------------------
% The truncated half-hook closure of the RMP O_L figure (arXiv:2011.12127
% Fig:MPDO-O_L): the ring seen edge-on, its return wire implied.  At each
% end of the traced row the wire runs half a stub out along the wire
% axis, turns through a 250-degree arc of radius stub/2 toward the return
% side, and truncates just short of the anchor line -- at 250 degrees the
% free tip sits cos(70) x r ~ 0.94 r past the arc's turning point, i.e.
% 0.03 x stub short of the line through the anchor, curling back toward
% the wire the way the paper's hooks do.  Reach bookkeeping: the hook
% protrudes stub/2 + r = one stub beyond the anchor, exactly the open-
% stub silhouette, and swings 2r = one stub toward the return side.
% Each hook is ONE path whose bend is ONE arc, so its rounding is uniform
% by construction (the plane_sweep3 defects: to[]-curves render too flat,
% multi-join racetracks mix sharp and rounded corners).  Both hooks bend
% toward the SAME return side -- where the racetrack's return wire would
% run -- and the arc angles pass through \tenkz_arcangle:n, so a vertical
% word's hooks wrap over the top and under the bottom toward the page's
% east margin, as in the reference figure.
\cs_new_protected:Npn \tenkz_draw_hooks:
  {
    \int_compare:nNnTF { \l__tenkz_nrows_int } = {1}
      { \tenkz_hook_row:nn {1}{below} }
      {
        \tenkz_hook_row:nn {1}{above}
        \tenkz_hook_row:nn {\l__tenkz_nrows_int}{below}
      }
  }

% hooks of row #1, bending toward side #2 (below|above, logical frame).
% Angle table (logical frame; the arc starts at the tip of the straight
% lead-out, whose position angle on the arc's circle is 90 for a bend
% below, 270 for a bend above; the sweep runs toward the anchor line,
% counter-clockwise exactly when the outward side sign and the bend sign
% agree):
%   below/west: 90 -> 340    below/east: 90 -> -160
%   above/west: 270 -> 20    above/east: 270 -> 520
\cs_new_protected:Npn \tenkz_hook_row:nn #1#2
  {
    % closures inherit (hue contract, clause 4): both hooks draw in the
    % traced row's resolved wire style
    \tenkz_row_wirestyle:nN {#1} \l__tenkz_cell_tl
    \tenkz_row_extent:n {#1}
    \int_compare:nNnF { \l__tenkz_first_int } = {0}
      {
        \tenkz_anchor:nnnN {#1}{\int_use:N \l__tenkz_first_int}
          {west} \l__tenkz_pb_tl
        \tenkz_anchor:nnnN {#1}{\int_use:N \l__tenkz_last_int}
          {east} \l__tenkz_pa_tl
        \str_if_eq:nnTF {#2} {below}
          {
            \tenkz_hook_end:Nnnn \l__tenkz_pb_tl {west} {90} {340}
            \tenkz_hook_end:Nnnn \l__tenkz_pa_tl {east} {90} {-160}
          }
          {
            \tenkz_hook_end:Nnnn \l__tenkz_pb_tl {west} {270} {20}
            \tenkz_hook_end:Nnnn \l__tenkz_pa_tl {east} {270} {520}
          }
        \tenkz@event{hooks|picture=\the\tenkz@pictureid|row=\int_eval:n{#1}|side=#2}
      }
  }

% one hook: anchor tl #1, outward side #2 (west|east, logical), logical
% start angle #3 and end angle #4.  The straight lead-out and the arc
% angles are frame-mapped; the arc's centre is implicit (tikz places it
% radius-away from the current point at the start angle), so the whole
% hook needs no absolute coordinates beyond the anchor.
\cs_new_protected:Npn \tenkz_hook_end:Nnnn #1#2#3#4
  {
    \draw[trace, \tl_use:N \l__tenkz_cell_tl]
      #1 -- ++\tenkz_side_vec:nn {#2} { \tenkz@dim{\tenkz@r@stub} / 2 }
      arc[start~angle=\tenkz_arcangle:n{#3},
          end~angle=\tenkz_arcangle:n{#4},
          radius=\dim_eval:n { \tenkz@dim{\tenkz@r@stub} / 2 }];
  }

% ---------- public entry points ----------
\NewDocumentEnvironment { tenkz } { O{} +b }
  { \tenkz_render:nn {#1}{#2} } { }
\NewDocumentCommand \tnpic { O{} m }
  { \tenkz_render:nn {#1}{#2} }

\ExplSyntaxOff
\endinput
